\documentclass[a4paper,11pt]{report}

% -------------------------------------
% PACKAGES
% -------------------------------------

% Fontenc: Western European punctuation
\usepackage[T1]{fontenc}

% UTF-8-Encoding
\usepackage[utf8]{inputenc}

% Language: English
\usepackage[english]{babel}

% Pictures
\usepackage{graphicx}

% Change font color
\usepackage[svgnames]{xcolor}

% Maths packages: 
\usepackage{amsmath,amssymb}

% Text symbols such as bullet or copyright:
\usepackage{textcomp}

% Allows an additional argument in an enumeration of the form \begin{enumerate}[EX i.]
\usepackage{enumerate}

% Quotations at the beginning of chapters
\usepackage{epigraph}

% Multicolumn environment
\usepackage{multicol}

% Appendix
\usepackage[toc, page]{appendix}

% Bibliography
\usepackage[
	backend=bibtex,
	style=nature,
	bibencoding=ascii
	%style=alphabetic
	%style=reading
]{biblatex}
\addbibresource{C:/Users/samue/Documents/Studium/Literatur/Mathematik-BA_Mathe.bib}

% Reference the last page by LastPage
\usepackage{lastpage}

% Custom theorem layouts: Theorems are defined within the document
\usepackage[
	thref, %theorem references by \thref
	amsmath, % Compatibility with amsmath
	thmmarks, % End marks
	hyperref % Compatibility with hyperref
]{ntheorem}


% References: should be loaded as the last package but before geometry
\usepackage{hyperref}
\hypersetup{
    bookmarksnumbered=true,
    bookmarksopen=false,
    bookmarksopenlevel=1,
    colorlinks=true,
    linkcolor=blue,
    urlcolor=DarkBlue,
    citecolor=DarkRed
}

% Page Layout
\usepackage{geometry}
\geometry{
	left=20mm,right=20mm,%
	bindingoffset=0mm, top=20mm,bottom=15mm
}

% -------------------------------------
% INFORMATION ON THE DOCUMENT
% -------------------------------------

\newcommand{\myauthor}{Samuel Lippl}
\newcommand{\mytitle}{A Mathematical Analysis of Backpropagation and Predictive Coding}
\newcommand{\shorttitle}{Mathematics of Predictive Coding}
\newcommand{\mysupervisor}{Prof. Dr. Markus Heydenreich}

% Custom headers and footers
\usepackage{fancyhdr}
\pagestyle{fancy}
\fancyhf{}
\fancyhead[L]{\shorttitle}
\fancyhead[C]{}
\fancyhead[R]{\myauthor}
\renewcommand{\headrulewidth}{0.4pt}
\fancyfoot[C]{\thepage}
\renewcommand{\footrulewidth}{0.4pt}

% -------------------------------------
% NEW COMMANDS
% -------------------------------------

\newcommand{\const}{\varsigma} % Denote a constant value
\newcommand{\var}{\chi} % Denote a variable
\newcommand{\conn}{\smile} % Denote a direct connection in either direction
\newcommand{\smallsquare}{\;\scalebox{0.6}{$\square$}\;}
\newcommand{\xor}{\;\dot{\vee}\;}
\newcommand{\In}{\text{Init}}
\newcommand{\St}{\text{St}}
\newcommand{\Val}{\text{Val}}
\newcommand{\tr}{\;\text{tr}}

\begin{document}

% -------------------------------------
% THEOREM COMMANDS
% -------------------------------------

\theoremstyle{change} % Number before Theorem
\theoremheaderfont{\normalfont\scshape}
\theorembodyfont{\normalfont}
\theoremseparator{}
\newtheorem{Par}{}[chapter]
\theoremseparator{.}
\newtheorem{The}[Par]{Theorem}
\newtheorem{Pro}[Par]{Proposition}
\newtheorem{Lem}[Par]{Lemma}
\newtheorem{Cor}[Par]{Corollary}
\newtheorem{Def}[Par]{Definition}
\newtheorem{Not}[Par]{Notation}
\newtheorem{Ex}[Par]{Example}
\newtheorem{Rem}[Par]{Remark}
\theoremstyle{nonumberplain}
\theoremheaderfont{\normalfont\itshape}
\theoremsymbol{$\square$}
\newtheorem{Bew}{Proof}

% -------------------------------------
% TITLEPAGE
% -------------------------------------

\begin{titlepage}
% Titlepage inspired by: https://www.overleaf.com/15991102jmmvbxxtbhfk#/61015645/

\newcommand{\HRule}{\rule{\linewidth}{0.5mm}} % Defines a new command for the horizontal lines, change thickness here

\center % Center everything on the page

%----------------------------------------------------------------------------------------
%	TITLE SECTION
%----------------------------------------------------------------------------------------

\HRule \\[0.4cm]
{ \huge \bfseries \mytitle}\\[0.4cm] % Title of your document
\HRule \\[1.5cm]

%----------------------------------------------------------------------------------------
%	HEADING SECTIONS
%----------------------------------------------------------------------------------------

\textsc{\LARGE Ludwig-Maximilians-Universität München}\\[0.5cm] % Name of your university/college
\textsc{\Large Fakultät für Mathematik, Informatik und Statistik}\\[1.5cm]


\textsc{\Large Bachelor of Mathematics}\\[0.5cm] % Major heading such as course name
\textsc{\large Bachelor's Thesis}\\[0.5cm] % Minor heading such as course title

%----------------------------------------------------------------------------------------
%	LOGO SECTION
%----------------------------------------------------------------------------------------

\vfill

\includegraphics[height=0.2\textheight]{Bilder/sigill.png}\\[1cm] % Include a department/university logo - this will require the graphicx package



%----------------------------------------------------------------------------------------
%	DATE SECTION
%----------------------------------------------------------------------------------------

\vfill

%----------------------------------------------------------------------------------------
%	AUTHOR SECTION
%----------------------------------------------------------------------------------------

{\Large June 07, 2018}\\[1cm]
\begin{Large}
\textsc{Author}\\Samuel Lippl\\[0.5cm]
\textsc{Supervisor}\\Prof. Dr. Markus Heydenreich
\end{Large}

\vfill
\end{titlepage}

\pagenumbering{roman}
\setcounter{page}{2}

% -------------------------------------
% SELBSTÄNDIGKEITSERKLÄRUNG (DEUTSCH ODER ENGLISCH?)
% -------------------------------------

\chapter*{Selbständigkeitserklärung}

Ich erkläre hiermit, dass ich die vorliegende Arbeit selbständig angefertigt, alle Zitate als solche kenntlich gemacht sowie alle benutzten Quellen und Hilfsmittel angegeben habe.

% -------------------------------------
% ABSTRACT
% -------------------------------------

\chapter*{Abstract}

% -------------------------------------
% INTRODUCTION
% -------------------------------------



% -------------------------------------
% TABLE OF CONTENTS
% -------------------------------------

\tableofcontents

% -------------------------------------
% INTRODUCTION
% -------------------------------------

\chapter{An Introduction to Computation in Machines, Brains and Mathematics}
\pagenumbering{arabic}
\setcounter{page}{1}

\epigraph{What kind of computer is the brain?}{Sejnowski et al. \cite[][1300]{Sejnowski1988}}

In "Science and Statistics" \cite{Box1976}, George E. P. Box discusses the proper statistical methodology of tackling scientific problems. Although not at all concerned with this discipline, the paper provides a remarkably clear perspective on the field of computational neuroscience, not only with respect to its methodology but also with respect to its findings. The article therefore provides a proper compass, to position this Bachelor's Thesis at the intersection of Machine Learning, computational neuroscience and, evidently, mathematics.\\
Computational neuroscience is concerned with explaining "how electrical and chemical signals are used to represent and process information" \cite[][1299]{Sejnowski1988}. These problems are tackled by constructing theoretical models that are then checked by empirical data. Like the methodology of many fields of science, this strategy corresponds to the general principle of "iteration between theory and practice" \cite[][791]{Box1976}. It is important that this loop is closed: while omission of theory leads to the unjustified application of a few recipes to problems that are not suitable for these analyses -- a problem which Box calls "cookbookery" \cite[][797]{Box1976} --, "mathematistry", the "development of theory for theory's sake", has the tendency "to redefine the problem rather than solve it" \cite[][797]{Box1976}.\\
Besides ensuring a valid analysis of the data, theory plays a second important role in science: making certain that the obtained explanations actually contribute to our understanding. Sejnowski et al. \cite{Sejnowski1988} situate neural models on a spectrum from realistic to simplifying:
The Hodgkin-Huxley model of signal transmission \cite{Hodgkin1952} in nerves is a perfect example of a realistic model. It describes the formation of an action potential at the level of a single neuron how an action potential, an electrical signal that is used for long-distance information transmission. Realistic models, i. e. those that attempt to describe the physiological properties of a neuronal network, are insufficient for two reasons:\\
Firstly, a huge amount of data is needed to fit these models, especially at the network level, and the computational power to simulate entire networks at this level is immense. However, even if the data were available -- which cannot be taken for granted, especially with concern to the human brain --, such a complex model may not enlighten our understanding of how the human brain works. A simplifying model may be better suited to illustrate important principles of information processing. These models are more concerned with how information is represented at a network level without considering the underlying physiological principles. Sejnowski et al. elaborate:
\begin{quote}
The study of simplifying models of the brain can provide a conceptual framework for isolating the basic computational problems and understanding the computational constraints that govern the design of the nervous system. \cite[][1300]{Sejnowski1988}
\end{quote}
This Bachelor's Thesis is concerned with the mathematical treatment of a popular simplifying model: neural networks. Keeping in mind the concerns that the previous arguments have raised, there are two balancing acts such a treatment has to accomplish:
\paragraph{Theory and Application}
Neural networks have originally been inspired by a biological model. Inherent to the concept of a simplifying model is, however, the possibility that such a network may become "an end in itself" \cite[][1305]{Sejnowski1988}. This outcome may be framed as a mathematization in Box's terms. However, this development has been integrated into the field of Machine Learning that is concerned with the question "How may we recognize patterns?" instead of the neuroscientific question "How may our brain recognize patterns?". Neural networks have become a powerful tool in Machine Learning in a way that cannot be reframed as a model of our brain. In particular, the backpropagation algorithm has been a popular method to perform gradient descent in neural networks that is not likely to be implemented by our brain. \cite{Liao2015} Other Machine Learning algorithms are even less related to a biological model of learning.\\
Does this mean that such a development is unsuitable for further neuroscientific discussion? No. The differences between how Machine Learning algorithms may be implemented in the brain and how learning in the brain occurs in reality provides fruitful ground for new ideas in both disciplines and discussion whether differences are due to constraints of evolution or optimality.\\
One example for such a mechanism is the development of the \emph{predictive coding model} that was first proposed and successfully applied by Rao and Ballard \cite{Rao1999} to explain certain observations in the visual cortex. Friston developed and generalized this model (see, for instance, \cite{Friston2005}). Recently, Whittington and Bogacz presented predictive coding in its relation to the backpropagation algorithm \cite{Whittington2017}.\\
It is helpful for these discussions to have a general framework for biologically plausible and implausible neural networks. I propose such a definition in chapter \ref{ch:theory}. It lays the groundwork for a mathematically rigorous and general introduction to backpropagation and predictive coding as well as their relations. This treatment may lead to predictive coding becoming just another end in itself. While empirical data is not of concern to a mathematical Bachelor's Thesis, I therefore introduce mathematical axioms in neural networks that have been supported by empirical data in order to prevent the mathematistry of my discussion of neural networks.
\paragraph{Generality and Conceptualization}
In the spirit of a mathematical analysis, I will introduce the concepts of prediction and learning in neural networks as well as backpropagation and predictive coding as generally as possible. This introduction shall be complemented by more restrictive conditions for which these networks will be simplified and easier to understand. An important focus of the Bachelor's Thesis will be the relation between the general framework and the more restrictive models.
\paragraph{}
Finally, this Bachelor's Thesis only presents the mathematical groundwork of some neural networks. Chapter \ref{ch:further-work} will present possible directions of future investigation and first endeavours within these directions.\\
These principles, I hope, will guide my attempt to present neural networks such that the analysis provides the rigorosity and generality of a mathematical treatment, is sufficiently grounded in reality to contribute to computational neuroscience and sufficiently concrete to provide clear results (or perspectives of clear results).

% -------------------------------------
% GROUNDS OF THEORY
% -------------------------------------

\chapter{Feedforward Neural Networks and Predictive Coding}\label{ch:theory}

This chapter will be concerned with introducing general neural networks. Section \ref{sec:theory-structure} will cover their structure. We will go on to states of neural networks in \ref{sec:theory-states} before introducing learning in neural networks in \ref{sec:theory-learning}. The specific theory of learning by backpropagation and predictive coding will be introduced alongside the more general groundwork.

\section{Structure of Neural Networks}\label{sec:theory-structure}

\begin{Par}[Neural Networks as Graphs]\label{theory-structure:intro}
Neural networks may be based on graphs where we refer to the nodes as \emph{Processing Units} (PU) and to the edges as \emph{Unit Connections} (UC). We can therefore use an adjacency matrix that specifies for each pair of PUs whether there is a connection between them (see \thref{par:adjacency}) to specify the neural network.\\
Moreover, there is a set of parameters that specify the behaviour of a neural network. As most of these parameters are directly connected to a UC, we will define a neural network by its \emph{weight matrix} which is a weighted adjacency matrix. However, we now face a difficulty concerning the relationship between the structure and the parameters. Generally, a connection between two PUs is indicated by a non-zero entry in the adjacency matrix while an entry of zero indicates that the two PUs are not connected. However, what if a connection between two PUs is weighted by zero? The difference between this case and no connection at all will become important in section \ref{sec:theory-learning}: a connection with weight zero may change during learning two unconnected PUs will always remain unconnected. Moreover, how may we define a neural network, where we know the UCs but not their weights?\\
Subsection \ref{subsec:theory-structure-matrix} presents a solution to these difficulties. Subsection \ref{subsec:theory-structure-forms} provides some vocabulary to talk about neural networks by introducing the standardized layerization. Subsequently, subsection \ref{subsec:theory-structure-predicode} will first introduce concepts from predictive coding, namely the generative model. This section will conclude with \ref{subsec:theory-structure-other} where we will embed the provided definition in the context of literature and will discuss the concept of \emph{parallelity} and its connection to neural networks.
\end{Par}

\subsection{The Network Matrix}\label{subsec:theory-structure-matrix}

\begin{Def}[$\const$-$\var$-extension of a set]
If $A$ is a set, we define the \emph{$\const$-$\var$-extension} 
\begin{equation}\label{eq:c-v-extension}
A_{\const,\var}:=\left(A\times\left\{\const,\var\right\}\right)\cup\{\var_A\},
\end{equation}
where we denote for any $a\in A$
\begin{equation}\label{eq:not-c-v}
a_{\const}\equiv(a,\const)\quad a_{\var}\equiv(a,\var)
\end{equation}
and may omit the subscripts when they become clear from context.\\
We will treat $0_{\const}$ like an ordinary zero, i. e. if a number in a matrix is omitted, it is $0_{\const}$ and a diagonal matrix only has $0_{\const}$ entries that are not on their diagonal.\\
If $a\in A\times\{\const\}$, then $a$ is \emph{constant} and we write $\St(a)=\const$ (as an abbreviation of \emph{status}. If $a\in\left(A\times\{\var\}\right)\cup\{\var\}$, then $a$ is \emph{variable} and we write $\St(a)=\var$. Finally, we denote the value of $a$ by $\Val(a_{\const})=\Val(a_{\var})=a$ where $\Val(\var)$ is not defined.
\end{Def}

\begin{Def}[$\const$-$\var$-extension of a function]
Consider sets $I, A^i$, where $i\in I$, $B$ and $C$.
\begin{enumerate}[a)]
\item
We define the \emph{passive} $\const$-$\var$-extension of a function $f:\prod_{i\in I}A^i\to B$ by
\begin{equation}\label{eq:domain-extension}
\begin{aligned}
f:\prod_{i\in I}A^i_{\const,\var}&\to B_{\const,\var},\\
\Val\left(f(a)\right)&:=\begin{cases}
f\left(\left(\Val(a_i)\right)_{i\in I}\right)&\text{ if }\neg\underset{i\in I}{\exists}a_i=\var,\\
c&\text{ if }\left(\underset{i\in I}{\exists}a_i=\var\right)\vee f\left(\left(\Val(a_i)\right)_{i\in I:a_i\ne\var},\cdot\right)\equiv c,\\
\var&\text{ else },
\end{cases}\\
\St\left(f(a)\right)&:=\begin{cases}
\const&\text{ if }\underset{i\in I}{\forall}\St(a_i)=\const,\\
\var&\text{ else }.
\end{cases}
\end{aligned}
\end{equation}
We will denote the passive extension by the same symbol as the original function.
\item
We define the \emph{active} $\const$-$\var$-extension of a function
\[
f:\prod_{i\in I}A^i\times C\to \prod_{i\in I}A^i\times B, 
\begin{pmatrix}
a\\c
\end{pmatrix}
\mapsto
\begin{pmatrix}
\left(f(a,c)_i\right)_{i\in I}\\
f_B(a,c)
\end{pmatrix}
\qquad
\]
in $A:=(A_i)_{i\in I}$ by
\begin{equation}
\begin{aligned}
f_{\const,\var:A}:\prod_{i\in I}A^i_{\const,\var}\times C&\to \prod_{i\in I}A^i_{\const,\var}\times B,\\
f{\const,\var:A}(a,c)_i:=
\begin{cases}
f(a,c)_i & \St(a_i)=\var\\
a_i & \St(a_i)=\const
\end{cases}.
\end{aligned}
\end{equation}
If $C=\emptyset$, we simply write $f_{\const,\var}$.
\end{enumerate}
\end{Def}

\begin{Par}[Passive Extension]
The passive extension allows application of a function $f$ to $\const$-$\var$-extended values. $f(a)$ is constant if and only if every argument is constant. Generally, the value of the extended $f$ is simply the value of the original $f$ when applied to the values of the arguments, for instance:
\[
1_{\const}+2_{\var}=(1+2)_{\var}=3_{\var}.
\]
This is covered by the first case of the definition. Things only get messy when the value of an argument is unknown. In this case, generally, $f(a)$ is unknown as well, for instance:
\[
1_{\const}+\chi=\chi.
\]
The only exception is given by a function that is constant in the unknown argument, for instance:
\[
0_{\const}\cdot\chi =0_{\var}.
\]
\end{Par}

\begin{Par}[Active Extension]
The active extension ensures that no constant value will be changed. This will, for instance, be useful in the training function, where we do not want $0_{\const}$ to change as this value encodes that there is no connection between two PUs, as we discussed in \thref{theory-structure:intro}.
\end{Par}

\begin{Par}[Adjacency Matrix]\label{par:adjacency}
An adjacency matrix has the following connection to a graph: Given a graph $\mathcal{G}:=(\mathcal{V},\mathcal{E}),\mathcal{E}\subseteq\mathcal{V}\times\mathcal{V}$, we may write $\mathcal{G}$ as a matrix $W\in\{0,1\}^{\mathcal{V}\times\mathcal{V}}$ where $W_{v,w}=1\Leftrightarrow (v,w)\in\mathcal{E}$.\\
If we consider weighted edges, we may generalize the notion to a matrix $W\in A^{\mathcal{V}\times\mathcal{V}}$ on a ring $A$ where $W_{v,w}\ne 0$ specifies the weight of edge $(v,w)\in\mathcal{E}$ and an edge with weight zero is synonymous with no edge at all. It is therefore important that the use of these weights is defined in an appropriate way.\\
We can further extend the notion of an adjacency matrix to $W\in A^{\mathcal{V}\times\mathcal{V}}_{\const,\var}$ where
\begin{equation}\label{eq:matrix-graph}
(v,w)\in\mathcal{E}\Leftrightarrow W_{v,w}\ne 0_{\const}.
\end{equation}
\end{Par}

\begin{Def}[Graph Relations]
If $\mathcal{G}:=(\mathcal{V},\mathcal{E})$ is a graph and $v,w\in\mathcal{V}$, we introduce the following three relations:
\begin{align}
v\to_{\mathcal{E}} w&\equiv (v,w)\in\mathcal{E},\\
v\leftarrow_{\mathcal{E}} w &\equiv w\to_{\mathcal{E}} v,\\
v\conn_{\mathcal{E}} w &\equiv v\to_{\mathcal{E}} w\vee v\leftarrow_{\mathcal{E}} w,
\end{align}
where we will often leave the subscript implicit.\\
We define their transitive closures $\to^*,\leftarrow^*,\conn^*$. Finally, if $v\to^* w$ (resp. $v\to w$) we say that $v$ precedes (resp. directly precedes) $w$, $w$ succeeds (resp. directly succeeds) $v$ and $v$ and $w$ are connected (resp. directly connected).\\
We will denote the set of directly preceding nodes by $\to v$ and the set of directly succeeding nodes by $v\to$.
\end{Def}

\begin{Def}[Connected Graph]
A graph is called \emph{connected} if and only if all its nodes are connected.
\end{Def}

\begin{Def}[Neural Network I: Topological Structure]
Consider a ring $A$. A tupel $\mathfrak{N}=(\mathcal{V},\theta), \theta\in A_{\const,\var}^{\mathcal{V}\times\mathcal{V}}$ is called a neural network if and only if the graph corresponding to $\theta$ is connected.
\end{Def}

\begin{Par}
The condition that neural networks have to be connected is not a restriction because if there are two unconnected parts in the network, we can simply consider these parts as separate networks.
\end{Par}

\begin{Not}[Adjacency Matrix]
While generally using the common matrix notations, we will refer to the elements of $\theta$ in a slightly different way. The entry in row $v$ and column $w$ is denoted by $\theta_{v\to w}$, the entries in row $v$ are denoted by $\theta_{v\to}$ and the entries in column $w$ are denoted by $\theta_{\to w}$. Finally, for $V,W\subseteq\mathcal{V}$
\[
\theta_{V\to W}\equiv\left(\theta_{v\to w}\right)_{v\in V,w\in W}.
\]
\end{Not}

\begin{Not}
For any family $(a_i)_{i\in I}$, if $J\subseteq I$, we denote $a_J\equiv(a_j)_{j\in J}$. If there is a set of functions $f_i:A_i\to B_i$, $i\in I$, and $a\equiv(a_i)_{i\in I},a_i\in A_i$, we denote $f(a)\equiv\left(f_i(a_i)\right)_{i\in I}$.
\end{Not}

\begin{Def}[Neural Network II: Functional Structure]
Let $\mathfrak{N}=(\mathcal{V},\theta)$ be a neural network.
Consider a set of \emph{global parameters} $\Lambda$, \emph{unit functions} $f\equiv\left(f^u\right)_{u\in\mathcal{V}}$ and \emph{integrative functions} $g\equiv\left(g^u\right)_{u\in\mathcal{V}}$.
The tupel $(\Lambda,f,g)$ is called the \emph{functional structure} of $\mathfrak{N}$ if and only if, for all $u\in\mathcal{V}$,
\begin{equation}
f^u:A\to A,\quad s^u\mapsto f^u(s^u)
\end{equation}
and
\begin{equation}
g^u:A^{\mathcal{V}}\times A_{\const,\var}^{\mathcal{V}\times\mathcal{V}}\times\Lambda\to A,
\quad
(s^{:o},\theta,\lambda)\mapsto g^u_{\theta,\lambda}(s^{:o})
\end{equation}
only depends on $s^{\to u:o}$, $\theta_{\to u}$ and $\lambda$.
\end{Def}

\begin{Par}[Integrative Functions]
A unit's integrative function summarizes the output of a unit's directly preceding nodes. The first argument therefore provides the values of these nodes, the scond argument provides the weight matrix $\theta$ and the third argument is a set of global parameters. How neural networks may take concrete states will be of concern in section \ref{sec:theory-states}.
\end{Par}

\begin{Par}[Unit Functions]
The unit functions relate the summarized input $s^{u:i}:=g^u_{\theta,\lambda}(s)$ to a node to its output value $s^{u:o}:=f^u(s^{u:i})$.
\end{Par}

\begin{Def}[Additive Neural Networks]
A neural network $\mathfrak{N}=(\mathcal{V},\theta)$ is called \emph{straight} if and only if
\[
\underset{v,w\in\mathcal{V}}{\forall}v\to w\Rightarrow\neg\underset{u\in\mathcal{V}}{\exists}v\to u\to^*w.
\]
A straight neural network $\mathfrak{N}=(\mathcal{V},\theta)$ with functional structure $(\Lambda,f,g)$ is called \emph{additive} if and only if, for all $u\in\mathcal{V}$,
\begin{equation}
g_{\theta}(s^{:o})=\theta^Ts^{:o}.
\end{equation}
$\mathfrak{N}$ is called \emph{general additive} neural network if and only if, for all $u\in\mathcal{V}$, there are sets $i_u,o_u\subseteq \mathcal{V}\setminus i_u$ such that
\begin{equation}
g^u_{\theta}(s^{i_u:o},s^{o_u:o})=\theta_{i_u\to u}^Ts^{i_u:i}+\theta_{o_u\to u}^Ts^{o_u:o}.
\end{equation}
\end{Def}

\begin{Par}
Straightness of neural networks ensures that there are no hidden dependencies. If $v$ is not only directly connected to $w$ but also indirectly, $g^w$ does not specify the entire dependence on the value of $v$. It is therefore crucial for an additive network to be straight; if not, the influence of the preceding nodes would seem additive according to $g$ but hidden dependencies may render it non-linear.
\end{Par}

\subsection{Layerization}\label{subsec:theory-structure-forms}

\begin{Par}[Layers of Neural Networks]
Many neural networks are defined in layers where $\mathcal{V}:=\mathcal{V}^{(0)}\dot{\cup}\dotsb\dot{\cup}\mathcal{V}^{(L)}$ and $v\to w\Leftrightarrow \underset{l\le L}{\exists}\left(v\in\mathcal{V}^{(l)}\wedge w\in\mathcal{V}^{(l+1)}\right)$. These layers can be extended to a class of neural networks called \emph{feedforward neural networks}. This subsection is concerned with finding and characterizing these layers, a process called \emph{layerization}.
\end{Par}

\begin{Not}
To simplify the notation of layers we define, for $l,l_1,\dotsc,l_n\in\{0,\dotsc,L\}$,
\[
V^{(\square l)}:=\bigcup_{k\in\{0,\dotsc,L\}:k\square l}V^{(k)},
\]
where $\square$ is some relation, e. g. $\square\equiv\le$, and
\[
V^{(l_1,\dotsc,l_n)}:=\bigcup_{i=1}^nV^{(l_i)}.
\]
\end{Not}

\begin{Def}[Feedforward Neural Networks]
A neural network is called feedforward if and only if its graph is acyclical. Feedforwardness is therefore a topological property.
\end{Def}

\begin{Pro}[Characterization of feedforward neural networks]\label{pro:chara-feedforward}
Let $\mathfrak{N}=(\mathcal{V},\theta)$ be a neural network and let $\theta$ be its weight matrix. Then, the following conditions are equivalent:
\begin{enumerate}[(i)]
\item
$\mathfrak{N}$ is feedforward,
\item
$\to$ is acyclical,
\item
There are layers $\mathcal{V}^{(0)},\dotsc,\mathcal{V}^{(L)}$ such that $\mathcal{V}=\mathcal{V}^{(0)}\dot{\cup}\dotsb\dot{\cup}\mathcal{V}^{(L)}$ and 
\begin{equation}
\mathcal{V}^{(l)}\ni v\to w\in\mathcal{V}^{(k)}\Rightarrow l<k,
\end{equation}
\item
There is a function $\pi:\mathcal{V}\to\left\{1,\dotsc,|\mathcal{V}|\right\}$ such that
\begin{equation}\label{eq:chara-feedforward-iv}
v\to w\Rightarrow \pi(v)<\pi(w),
\end{equation}
\item
There is a function $\pi:\mathcal{V}\to\left\{1,\dotsc,|\mathcal{V}|\right\}$ such that
\begin{equation}\label{eq:chara-feedforward-v}
\theta_{\pi}:=\left(\theta_{v\to w}\right)_{\pi(v),\pi(w)}
\end{equation}
is an upper triagonal matrix.
\end{enumerate}
\end{Pro}

\begin{Bew}
\textit{(i)$\Leftrightarrow$(ii)}
Follows by definition.\\
\textit{(i)$\Rightarrow$(iii)}
is a special case of \thref{cor:standardized-layerization-feedforward}.\\
\textit{(iii)$\Rightarrow$(iv)}
Assign numbers in $\{1,\dotsc,|\mathcal{V}^{(0)}|\}$ to the elements of $\mathcal{V}^{(0)}$. For any $0<l\le L$ assign numbers in $\left\{\sum_{j=1}^{l-1}|\mathcal{V}^{(j)}|+1,\dotsc,\sum_{j=1}^{l}|\mathcal{V}^{(j)}|\right\}$ to the elements of $\mathcal{V}^{(l)}$. Then $v\to w$ where $v\in\mathcal{V}^{(l)}$ and $w\in\mathcal{V}^{(k)}$ implies $l<k$ and therefore 
\[
\pi(v)\le \sum_{j=1}^{l}|\mathcal{V}^{(j)}|<\sum_{j=1}^{l}|\mathcal{V}^{(j)}|+1\le \sum_{j=1}^{k}|\mathcal{V}^{(j)}|+1\le\pi(w).
\]
\textit{(iv)$\Rightarrow$(v)}
We simply prove that (\ref{eq:chara-feedforward-iv}) implies (\ref{eq:chara-feedforward-v}). This becomes clear from the fact that 
\[
(\theta_{\pi})_{\pi(v),\pi(w)}\ne0_{\const}\Rightarrow\theta_{v\to w}\ne0_{\const}\Rightarrow v\to w\Rightarrow\pi(v)<\pi(w)
\]
which is the definition of an upper triagonal matrix.\\
\textit{(v)$\Rightarrow$(iv)}
We now prove that (\ref{eq:chara-feedforward-v}) implies (\ref{eq:chara-feedforward-iv}). This follows from
\[
v\to w\Rightarrow(\theta_{\pi})_{\pi(v),\pi(w)}\ne0_{\const}\Rightarrow\pi(v)<\pi(w).
\]
\textit{(iv)$\Rightarrow$(ii)}
Clearly, (\ref{eq:chara-feedforward-iv}) could not hold if $\to$ contained a cycle.
\end{Bew}

\begin{Par}
Condition (iii) has a more universal appeal which we will discover now. These layers can be specified either by a partition of $\mathcal{V}$ or a transitive total order $\preceq$. After proving their equivalence, we go on to provide the definition of different kinds of layerizations.
\end{Par}

\begin{Def}[Layerization]\label{def:layerization}
A transitive total (but not necessarily antisymmetric) order $\preceq$ on $\mathcal{V}$ is called layerization. We denote
\[
v\prec w\equiv v\preceq w\wedge w\npreceq v\qquad v\simeq w\equiv v\preceq w\wedge w\preceq v.
\]
A connection $v\to w$ is called
\begin{itemize}
\item
\emph{forward} if and only if $v\prec w$,
\item
\emph{lateral} if and only if $v\simeq w$,
\item
\emph{backward} if and only if $v\succ w$.
\end{itemize}
The corresponding equivalence classes $\mathcal{V}/{\sim}$ are called \emph{layers}. As $\preceq$ is transitive, we may define $\preceq$ on $\mathcal{V}/{\simeq}$ by
\begin{equation}
V\preceq W\equiv\underset{v\in V, w\in W}{\forall}v\preceq w\equiv\underset{v\in V,w\in W}{\exists}v\preceq w.
\end{equation}
We can now notate the layerization as a partition $\mathcal{V}=\dot{\bigcup}_{l=0}^L\mathcal{V}^{(l)}$ where
\[
\mathcal{V}/{\simeq}=\left\{\mathcal{V}^{(l)}:0\le l\le L\right\}
\qquad 
\mathcal{V}^{(l)}\preceq\mathcal{V}^{(l+1)},
\]
which yields the layerization form of \thref{pro:chara-feedforward}.\\
We will call $\preceq$ the \emph{ordinal} and the partition $\left(\mathcal{V}^{(l)}\right)_{0\le l\le L}$ the \emph{set notation} of the layerization. $L$ is called the \emph{depth} of the layerization.
\end{Def}

\begin{Pro}\label{pro:layerization-11}
There is a one-to-one correspondence between ordinal and set layerization.
\end{Pro}

\begin{Bew}
The function from ordinal to set layerization has already been given by \thref{def:layerization}. On the other hand, if there is a set layerization, the order
\[
\mathcal{V}^{(l)}\ni v\preceq w\in\mathcal{V}^{(k)}\equiv l\le k
\]
is total and transitive.\\
We now only have to prove that the ordinal notation converted to the set notation converted to the ordinal notation provides indeed the same layerization. Let $v,w\in\mathcal{V}$. WLOG, suppose that $v\preceq w$. Then $v\prec w\xor v\simeq w$ where $\xor$ is the exclusive disjunction. If $v\prec w$, there are different equivalence classes $V\ne W$ such that $v\in V,w\in W$. As $v\prec w$, $V\prec W$ and therefore layer $V$ comes before layer $W$. This, again, yields $v\prec w$. The other case, $v\simeq w$, works in an analogous way.
\end{Bew}

\begin{Par}
By \thref{pro:layerization-11}, we will use the relational and set form of layerization interchangeably from now on. The next definition provides an important stepping stone towards the function of layerizations by allowing us to differentiate a connection $v\conn w$ not only into direction of information but also the global network direction: the flow\footnote{
which is different from the standard flow in graph theory as defined, for instance, in \cite[][149]{Diestel2017}}. The main motivation behind this is \thref{def:layerization-equivalence}.
\end{Par}

\begin{Def}
Given a layerization $\preceq$ of a network $\mathcal{N}$, we define the following relations which we will refer to as the \emph{flow} of $\mathcal{N}$.
\begin{equation}
\begin{aligned}
v\trianglerighteq w&\equiv 
v\conn w\wedge v\preceq w,&
v\trianglelefteq w&\equiv w\trianglerighteq v,&
v\smallsquare w&\equiv v\trianglerighteq w\wedge v\trianglelefteq w,\\
v\triangleright w &\equiv v\trianglerighteq w\wedge\neg(v\smallsquare w),&&&
v\triangleleft w&\equiv w\triangleright v.
\end{aligned}
\end{equation}
\end{Def}

\begin{Lem}
$\underset{v,w\in\mathcal{V}}{\forall}v\conn w\Leftrightarrow v\smallsquare w\xor v\triangleright w\xor v\triangleleft w$ where $\xor$ is an exclusive disjunction.
\end{Lem}

\begin{Bew}
Suppose $\neg(v\conn w)$. It is evident from its definition that there is no flow between $v$ and $w$ in this case. In turn, if $v\conn w$, the totality of $\preceq$ implies the trichotomy and the definition of the strict inequalities implies that the trichotomy is exclusive.
\end{Bew}

\begin{Def}
A layerization $\preceq$ on $\mathcal{V}$ is called
\begin{enumerate}[a)]
\item
\emph{minimal} if and only if there is a sequence $v_0,\dotsc,v_L$ such that $v_l\in\mathcal{V}^{(l)}$ and $v_l\triangleright v_{l+1}$ for all $l=0,\dotsc, L-1$,
\item
\emph{spanned} if and only if the minimal and maximal values of $\preceq$ and $\triangleright^*$, the transitive closure of $\triangleright$, are identical, 
\item
\emph{leftwards} (resp. \emph{rightwards}) \emph{anchored} if and only if for any $v_l\in\mathcal{V}^{(l)}$ there is a sequence $v_0\triangleright\dotsb\triangleright v_l$ (resp. $v_l\triangleright\dotsb\triangleright v_L$).
\end{enumerate}
\end{Def}

\begin{Lem}
An anchored layerization is minimal.
\end{Lem}

\begin{Bew}
If the layerization is leftwards anchored, set $l=L$. If the layerization is rightwards anchored, set $l=0$. As the set form is a partition, $\mathcal{V}^{(l)}\ne\emptyset$ and minimality follows.
\end{Bew}

\begin{Def}\label{def:layerization-equivalence}
Two layerizations $\preceq_1$, $\preceq_2$ of a neural network $\mathfrak{N}$ are called \emph{equivalent} (Not.: $\preceq_1\sim\preceq_2$) if and only if the flow remains identical.
\end{Def}

\begin{Par}
Layerizations provide different perspectives on a neural network. Generally, however, we want our notion of forward, lateral or backward connections (especially whether a given network is feed-forward) to remain invariant. Equivalent layerizations ensure such invariance.
\end{Par}

\begin{Lem}\label{lem:depth}
Given a layerization $\preceq$ and a chain $v_0\triangleright\dotsb \triangleright v_n$ where $v_0\in\mathcal{V}^{(l)}$ and $v_n\in\mathcal{V}^{(k)}$, we obtain $k-l\ge n$.
\end{Lem}

\begin{Bew}
Proof by induction on $k-l$. By definition, $v_n$ has to lie in a higher layer than $v_0$ if $v_0\ne v_n$. Therefore, if $k<l$, such a chain is not pssoible and if $k=l$, $n=0$.\\
If $k-l\ge 1$, consider $v_{n-1}\in\mathcal{V}{(k')}$ and observe that $k'<k$ and therefore $k'\le k-1$. Then, by induction hypothesis, $n-1\le k'-l\le k-1-l$ and therefore $n\le k-l$. 
\end{Bew}

\begin{Pro}[Depth of Layerizations]
Given a layerization $\preceq$, any equivalent minimal layerization has the same depth $L$. Moreover, every equivalent layerization that is not minimal has a depth of at least $L+1$. We call $L$ the depth of $\preceq$.
\end{Pro}

\begin{Bew}
By \thref{lem:depth}, any layerization must at least have the depth of the longest chain $v_0\triangleright\dotsb\triangleright v_L$. As $\triangleright$ is invariant under equivalent layerizations, all equivalent minimal layerizations have the same depth $L$. On the other hand, if a layerization has the depth $L$, any chain $v_0\triangleright\dotsb\triangleright v_L$ suffices the condition of minimality because every $v_l$ must be in a distinct layer and if $v_l\triangleright v_k$, $v_l$ must be in a lower layer than $v_k$.
\end{Bew}

\begin{The}[Standardized Layerization]\label{the:standardized-layerization}
For every flow $\triangleright$ of a neural network $\mathcal{N}$, there is a unique minimal spanned and leftwards (resp. rightwards) anchored layerization $\preceq_l$ (resp. $\preceq_r$) such that the flow of $\preceq_l$ (resp. $\preceq_r$) is $\triangleright$. We call $\preceq_l$ (resp. $\preceq_r$ \emph{leftwards} (resp. \emph{rightwards}) \emph{standardized layerization}.
\end{The}

\begin{Bew}
If there are no $v,w\in\mathcal{V}$ such that $v\triangleright w$, a minimal layerization must have depth $L=0$ and we therefore obtain $\mathcal{V}^{(0)}=\mathcal{V}$ for both the leftwards and the rightwards standardized layerization. This is equivalent to the original layerization: if $v\smallsquare w$, then $v\conn w$ and therefore $v\smallsquare_lw$ and $v\smallsquare_rw$. This proves equivalence as there are no $v,w$ such that $v\triangleright w$. The minimal and maximal values of $\preceq_l$, $\preceq_r$ and $\triangleright^*$ are given by $\mathcal{V}^{(0)}=\mathcal{V}^{(L)}=\mathcal{V}$ and therefore, $\preceq_l$ and $\preceq_r$ are spanned. Anchoredness is rendered trivial which proves the first case. We will therefore now suppose
\begin{equation}\label{eq:standardized-layerization-1}
\underset{v,w\in\mathcal{V}}{\exists}v\triangleright w
\end{equation}
and define $\mathcal{V}^{(0)}$ as the minimal values of $\triangleright^*$ and $\mathcal{V}^{(\max)}$ as the maximal values of $\triangleright^*$ which is necessary as the layerization must be spanned. These two sets are disjunct: suppose there is some $u\in\mathcal{V}^{(0)}\cap\mathcal{V}^{(L)}$. By definition, $u$ must be connected to $v$ from (\ref{eq:standardized-layerization-1}), i. e. there is some chain $u\conn v_1\conn\dotsb\conn v_{n-1}\conn v$. We prove $u\smallsquare^*v$ by induction on $n$:\\
If $n=1$, $u\conn v$ and therefore $u\triangleright v\vee u\smallsquare v\vee u\triangleleft v$. As $u$ is both a maximal and a minimal element, this implies $u\smallsquare v$ and therefore $u\smallsquare^* v$.\\
Now suppose $n>1$. By induction hypothesis, we know that $u\smallsquare^* v_{n-1}$. On the other hand, as $v_{n-1}\conn v$. As $\trianglerighteq^*$ is transitive, $v_{n-1}$ is therefore also maximal. By the same argument as above, $v_{n-1}\smallsquare^* v$ and therefore $u\smallsquare^* v$.\\
The same method yields $u\smallsquare^* w$ and therefore $v\smallsquare^* w$ which is a contradiction with (\ref{eq:standardized-layerization-1}) and proves that $\mathcal{V}^{(0)}\cap\mathcal{V}^{(L)}=\emptyset$.\\
We apply the following algorithm to obtain the remaining layers if we want a leftwards anchored layerization:
\begin{enumerate}
\item
Set $t=0$ and $R^{(0)}:=\mathcal{V}\setminus\left(\mathcal{V}^{(0)}\cup\mathcal{V}^{(L)}\right)$,
\item
While $R^{(t)}\ne\emptyset$:
\begin{enumerate}
\item
Define $\mathcal{V}^{(t+1)}$ as the set of minimal values in $R^{(t)}$,
\item
Set $R^{(t+1)}:=\mathcal{R}^{(t)}\setminus\mathcal{V}^{(t+1)}$,
\item
Increment $t$ by $1$,
\end{enumerate}
\item
Set $L:= t$.
\end{enumerate}
Clearly, this yields a partition. As $R^{(L)}$ is empty, every $v\in\mathcal{V}$ is in some layer and the construction ensures that the sets $\mathcal{V}^{(0)},\dotsc,\mathcal{V}^{(L-1)}$ as well as the sets $\mathcal{V}^{(1)},\dotsc,\mathcal{V}^{(L)}$ where we have already proven that $\mathcal{V}^{(0)}\cap\mathcal{V}^{(L)}=\emptyset$.\\
The layerization is leftwards anchored as we prove by induction on $l$ in the definition: If $l=0$, the criterion is trivial. If $l>0$, as $v_l$ was minimal in $R^{(l-1)}$ but not in $R^{(l-2)}$ there must be some $v_{l-1}\in\mathcal{V}^{(l-1)}$ such that $v_{l-1}\triangleright v_l$ and the chain follows by induction hypothesis.\\
On the other hand, the layers $\mathcal{V}^{(1)},\dotsc,\mathcal{V}^{(L-1)}$ are uniquely defined by the condition that the layerization must be leftwards anchored. Take any $v\in\mathcal{V}^{(l)}$, $1\le l\le L-1$, and suppose that there is an equivalent, leftwards anchored and spanned layerization $\preceq'$ where $v\in\mathcal{V}'^{(k)},k\ne l$. If $k<l$, $\preceq'$ would be impossible by \thref{lem:depth} as there is a chain $v_0\triangleright\dotsb\triangleright v_l$ as $\preceq_l$ is leftwards spanned. If $k>l$, $\preceq_l$ would be impossible as $\preceq'$ is leftwards spanned but we have already shown that it is well-defined. There $k=l$ and $\preceq_l$ is uniquely defined.\\
An analogous method shows that $\preceq_r$ is unique and well-defined where we use a top-down instead of a bottom-up approach with the algorithm and all arguments thereafter.
\end{Bew}

\begin{Cor}
For every layerization $\preceq$ on a neural network $\mathfrak{N}$, there is a unique equivalent leftwards or rightwards standardized layerization.
\end{Cor}

\begin{Bew}
The layerization is provided by the flow defined by $\preceq$ and \thref{the:standardized-layerization}.
\end{Bew}

\begin{Cor}\label{cor:standardized-layerization-feedforward}
If $\mathfrak{N}$ is a feed-forward network and $v\trianglerighteq w\equiv v\to w$, there is a unique leftwards or rightwards standardized layerization.
\end{Cor}

\begin{Par}
We have therefore found a way to layerize any feed-forward network in a sensible way. On the other hand, when we extend the feed-forward network to a more general network, we can still talk about the network's layers in a standardized manner. The next section provides the groundwork for this very necessity.
\end{Par}

\subsection{The Generative Model}\label{subsec:theory-structure-predicode}

\begin{Def}[Stochastic Networks]
Consider a probability space $(\Omega,\mathcal{F},\mathbb{P})$. We define a \emph{stochastic network} $\mathfrak{S}:=(\mathcal{V},\theta)$ with functional structure $(\Lambda,f,g)$, $\Lambda=\Omega\times\Lambda'$ such that, for all $u\in\mathcal{V}$,
\begin{equation}
g^u_{\theta,(\omega,\lambda')}\left(s\right)=g^u_{\theta,\lambda'}(\omega)
\end{equation}
where
\[
g^u_{\theta,\lambda'}:\Omega\to\mathbb{R}
\]
is a random variable. We may then consider the entire network as a random vector
\[
g_{\theta,\lambda'}:\Omega\to\mathbb{R}^{\mathcal{V}}
\]
\end{Def}

\begin{Par}[Predictive Coding I: The Generative Model]\label{predicode-1}
The notion of stochastic networks allows us to lay the groundwork for predictive coding: the neural model of reality which is, according to the theory, given by a \emph{hierarchical model}: a feed-forward network $\mathfrak{H}=(\mathcal{V},\theta)$ where the distribution of the lower-level PUs is specified in dependence on the realization of their preceding, higher-level PUs. This is motivated in an applied context by characterizing the higher-level PUs as causes of the lower-level PUs. Generally, this means that $\to v$ is the set of causes of $v$. In turn, $v$ is a cause for all $v\to$.\\
As an example, we will consider a small sector of a generative model that is concerned with vision. It is of interest to many organisms to estimate the speed of an object being thrown towards them. If we restrict ourselves to visual information, the only knowledge that is available to the organism is the input to the retinal receptors
\[
\text{ret}\subset\mathcal{V}.
\]
The retinal input is modeled as stochastically dependent on the properties of the object being thrown. For instance, there may be a PU
\[
\text{ball}\in\mathcal{V}
\]
encoding whether the object being thrown is a ball. As the information that an object is being thrown is not a part of the network, the value of that PU is not dependent on any other PU. We thus have an integrative function representing the \emph{a priori} distribution of \text{ball}:
\[
g^{\text{ball}}_{\theta,\lambda'}\sim\text{Bernoulli}(p_{\lambda'}),
\quad
0\le p_{\lambda'}\le 1.
\]
(As $\text{ball}$ is a highest-level PU, $p$ cannot depend on $\theta$.)\\
It is possible that the model does not contain any information on the likelihood of a given object being a ball. In this case, $\text{ball}$ would not be part of the actual stochastic network but rather a value that has to be specified beforehand and determines the distribution of the stochastic network itself. We will not consider this case as it is less realistic than one might think at first (see \ref{par:missing-priors}).\\
The speed of an object may now depend on whether this object is a ball, i. e., if
\[
\text{speed}\in\mathcal{V}
\]
encodes the speed of the objects, we may observe that 
\[
\text{ball}\to\text{speed}.
\]
The speed itself may be normally distributed\footnote{
Negative speed would encode the object moving in the reverse direction.} given its causes:
\[
g^{\text{speed}}_{\theta,\lambda'}(s)\sim\mathcal{N}\left(\mu^{\text{speed}}_{\theta,\lambda'}(s),\Sigma^{\text{speed}}_{\theta,\lambda'}(s)\right),
\quad
\mu^{\text{speed}}_{\theta,\lambda'},\Sigma^{\text{speed}}_{\theta,\lambda'}:\mathbb{R}^{\mathcal{V}}\to\mathbb{R}.
\]
where $\mu^{\text{speed}}$ and $\Sigma^{\text{speed}}$ are parameters that specify the normal distribution. More generally, we may have some (increasing) distribution
\[
p^{\text{speed}}_{\theta,\lambda'}(s):\mathbb{R}\to [0,1]
\]
of $g^{\text{speed}}_{\theta,\lambda'}(s)$.\\
Similarly, the speed of the objects may shape the retinal input:
\[
\text{speed}\to\text{ret}.
\]
Under comparable assumptions (but keeping in mind that there are multiple retinal inputs), we may conclude
\[
g^{\text{ret}}_{\theta,\lambda'}(s)\sim\mathcal{N}\left(\mu^{\text{ret}}_{\theta,\lambda'}(s),\Sigma^{\text{ret}}_{\theta,\lambda'}(s)\right),
\quad
\mu^{\text{speed}}_{\theta,\lambda'}:\mathbb{R}^{\mathcal{V}}\to\mathbb{R}^{\text{ret}},
\quad
\Sigma^{\text{speed}}_{\theta,\lambda'}:\mathbb{R}^{\mathcal{V}}\to\mathbb{R}^{\text{ret}\times\text{ret}}.
\]
Again, we may replace these parameters with a general multivariate distribution
\[
p^{\text{ret}}_{\theta,\lambda'}(s):\mathbb{R}^{\text{ret}}\to[0,1].
\]
We have thus provided a path from a higher-level cause towards the input onto the retina:
\[
\text{ball}\to\text{speed}\to\text{ret}
\]
In general, however, the higher-level cause is not only not known, it is not even empirically measurable. How we infer higher-level causes from lower-level input and how we construct these causes will be the concern of the following section \ref{sec:theory-states}. For now, we will provide a general definition of a generative model.
\end{Par}

\begin{Not}
We denote independence between random variables $X$ and $Y$ by
\[
X\bot Y.
\]
\end{Not}

\begin{Def}
Consider a set $X$ and a probability space $(\Omega,\mathcal{F},\mathbb{P})$. Stochastic functions $f,g:X\times\Omega$ are called \emph{independent} if and only if
\[
\underset{x\in X}{\forall}f(x)\bot g(x).
\]
\end{Def}

\begin{Def}[Generative Models]\label{def:generative-model}
Consider a feed-forward stochastic network $\mathfrak{G}=\left(\mathcal{G},\xi\right)$ with functional structure $(\Lambda'\times\Omega,f,g)$ and distribution
\[
g_{\xi,\lambda'}(s)\sim\mathbb{P}_{\xi,\lambda'}(s).
\]
$\mathbb{P}_{\xi,\lambda'}$ is called \emph{generative distribution} if and only if
\[
\underset{v,w\in\mathcal{G}}{\forall}v\to^* w\Rightarrow g^v\bot g^w.
\]
In this case $\mathfrak{G}$ is called \emph{generative model}.
\end{Def}

\begin{Def}[Normal Generative Model]
A generative model $\mathfrak{G}$ is called \emph{normal generative model} if and only if there is a non-stochastic functional structure $(\Lambda',f,\mu)$ and, for any $\lambda'\in\Lambda'$, some covariance matrix
\[
\Sigma_{\lambda'}\in\mathbb{R}^{\mathcal{G}\times\mathcal{G}}
\]
such that
\[
g_{\xi,\lambda'}(s)=\mu_{\xi,\lambda'}(s)+\epsilon_{\lambda'},
\quad
\epsilon_{\lambda'}\sim\mathcal{N}(0,\Sigma_{\lambda'}).
\]
\end{Def}

\begin{Par}
Because a normal generative model is also a generative model, we know that
\[
\underset{v,w\in\mathcal{G}}{\forall}v\to^* w\Rightarrow\Sigma_{v\to w}=0.
\]
\end{Par}

\subsection{Alternative Definitions of Neural Networks}\label{subsec:theory-structure-other}

\begin{Par}
This subsection is dedicated to a brief embedding of the provided definition of neural networks into the literature. I have mostly based my definition on \cite{Rojas1996}. However, my definition is more explicit than the introduction and more general than the later sections, mostly because a mathematical treatment tries to differentiate between applications and the general definitions.
\end{Par}

\begin{Par}[Parallelity]
Parallelity is a crucial aspect of neural networks. Massively parallel pathways allows the transformation of input space that is necessary for tasks like visual recognition. On the other hand, parallel pathways reduce the error rate of unreliable neural networks. These considerations have led Guresen and Kayakutlu \cite{Guresen2011} to claim parallelity as a necessary feature of a neural network. In their proposed definition, they require that "at least two of the multiple [Processing Units are] connected in parallel" \cite[][428]{Guresen2011} because, they agure, "structures [...] with one or more [Processing Units] connected serially cannot be referred as [a neural network] because it will lose the power of parallel computing and starts to act more like existing computers than a brain" \cite[][428]{Guresen2011}.\\
In principle, I agree with their point. I would argue, however, that it is not helpful to include such a criterion in the definition. After all, if there is one parallel pathway in a massive serial network, this would correspond to a neural network according to their notion even though the difference to a uniquely serial network is insubstantial.\\
The definition therefore seems to raise an issue that cannot be resolved and is analogous to the definition of a "heap of sand", another instance of the Sorites paradox (see \cite{Hyde2014}): without doubt, one grain of sand is necessary to have a heap but it is not sufficient and we can probably not find a hard definition of such a vague statement. We would be better advised to simply define the underlying property a "heap" refers to, i. e. the \emph{number} of grains of sand that allows us to use more rigorous notions.\\
Comparably, parallel processing such that it is advantageous cannot be covered by such a broad definition -- not only does parallelity exist on a spectrum, it also strongly depends on the context. I would therefore suggest that an attempt to find a way to talk about parallelity, for instance by a \emph{degree of parallelity} may be more fruitful than Guresen and Kayakutlu's hard borders. It is, after all, rather unlikely that exclusion of a few pathological cases would improve a mathematical definition.
\end{Par}

\section{States of Neural Networks}\label{sec:theory-states}

\begin{Par}
In the previous section, we have defined the topology of neural networks: their connectionist pattern and their functional structure. It remains an open question how this topology determines the values of the neural network. We will therefore formally introduce the notion of \emph{states} that has already been used in the previous section. \emph{State functions} will provide a general mechanism to relate states of input PUs with states of the entire network. The \emph{canonical state functions} provide an intuitive connection between the functional structure of a network and its state attribution and are the ultimate goal of subsection \ref{subsec:theory-states-basics}.\\
These are especially clear if we consider feed-forward neural networks, where the input is provided to the lowest layer $\mathcal{V}^{(0)}$. On the other hand, the notation is especially simple if the network is additive. We will therefore explore these cases in subsection \ref{subsec:theory-states-add}.\\
Finally, the network functions provide an appropriate framework for inference in predictive coding. We will introduce the concept of \emph{free energy} and then introduce inference in subsection \ref{subsec:theory-states-predicode}.
\end{Par}

\subsection{Network Functions}\label{subsec:theory-states-basics}

\begin{Def}[States]
Given a neural network $\mathfrak{N}=(\mathcal{V},\theta)$, we will call
\[
s\in S[\mathfrak{N}]:=A^{\{i,o\}\times\mathcal{V}}
\]
a \emph{state} of $\mathfrak{N}$. If $(k,v)\in\{i,o\}\times\mathcal{V}$, we denote the $(k,v)$-entry in $s$ by $s^{v:k}$. We refer to $s^{v:i}$ as the \emph{input value} of $v$ and $s^{v:o}$ as the \emph{output value} of $v$.\\
Furthermore, we denote, for $v\in\mathcal{V},V\subseteq\mathcal{V}$ and $k\in\{i,o\}$
\[
s^{V:k}\equiv\left(s^{v:k}\right)_{v\in V},
\quad
s^{:k}\equiv s^{\mathcal{V}:k},
\quad
s^{v:}\equiv\left(s^{v:i},s^{v:o}\right),
\quad
s^{V:}\equiv\left(s^{v:}\right)_{v\in V}.
\]
We define the \emph{$\const$-$\var$-extended state} $s\in S_{\const,\var}[\mathfrak{N}]:=A_{\const,\var}^{\{i,o\}\times\mathcal{V}}$.\\
Using the analogous notation, we define the \emph{state} of $\mathfrak{N}$ in $V\subseteq\mathcal{V}$ by 
\[
S^V[\mathfrak{N}]:=A^{\{i,o\}\times V}
\]
and the \emph{$\const$-$\var$-extended state} by
\[
S^V_{\const,\var}[\mathfrak{N}]:=A^{\{i,o\}\times V}_{\const,\var}.
\]
\end{Def}

\begin{Par}[Input and Output Values]
The input value of a PU $v$ $s^{v:i}$ refers to the value that is fed into a particular node whereas the output value of a PU refers to the value it distributes to its succeeding nodes. Details can be found below in \thref{def:balancedness}.
\end{Par}

\begin{Par}[Extended States]
An extended state allows notation of invariant states that are not allowed to change. On the other hand, we can specify that the states $\var$ are unknown.
\end{Par}

\begin{Def}[State Functions]
Consider a neural network $\mathfrak{N}=(\mathcal{V},\theta)$. We define the \emph{input} $\mathcal{I}\subseteq\mathcal{V}$ and the \emph{output} $\mathcal{O}\subseteq\mathcal{V}$. Then, 
\begin{equation}\label{eq:state-function}
T:S_{\const,\var}^{\mathcal{I}:i}[\mathfrak{N}]\times A_{\const,\var}^{\mathcal{V}\times\mathcal{V}}\times\Lambda\to S^{\mathcal{O}}_{\const,\var}[\mathfrak{N}],\qquad 
(i,\theta,\lambda)\mapsto T_{\theta,\lambda}(i)
\end{equation}
a state function and consider the following special cases:
\begin{enumerate}[a)]
\item
If $\mathcal{O}=\mathcal{V}$, $T$ is called a \emph{network function}.
\item
If $\mathcal{I}=\mathcal{O}=\mathcal{V}$, $T$ is called a \emph{step}.
\item
A network function $\In_a$ such that $\In_a^{\mathcal{I}}(i)=i$ and $\In_a^{\mathcal{V}\setminus\mathcal{I}}(i)=\left(a\right)_{v\in\mathcal{V}\setminus\mathcal{I}}$ where $a\in A_{\const,\var}$ is called $a$-initialization. Where this is clear from context, we generally identify any input vector $i$ with its $\var$-initialization $\In\equiv\In_{\var}$.
\end{enumerate}
\end{Def}

\begin{Not}[Defining state functions]
It is rather bothersome to define extended functions. We may therefore define the state function by 
\[
\tilde{T}:S^{\mathcal{I}:i}[\mathfrak{N}]\times A^{\mathcal{V}\times\mathcal{V}}\times\Lambda\to S^{\mathcal{O}}[\mathfrak{N}]
\]
and then define the active extension in $S^{\mathcal{I}\cap\mathcal{O}}[\mathfrak{N}]$ as $T$ (along with the implicit passive extensions), i. e.
\begin{equation}
T\equiv\tilde{T}_{\const,\var:S^{\mathcal{I}\cap\mathcal{O}:i}[\mathfrak{N}]}.
\end{equation}
On the other hand, we may also define state functions by
\[
\tilde{T}:S^{\mathcal{I}}[\mathfrak{N}]\times A^{\mathcal{V}\times\mathcal{V}}\times\Lambda\to S^{\mathcal{O}}[\mathfrak{N}],
\]
where the output value of the state $s\in S^{\mathcal{I}}[\mathfrak{N}]$ has no influence on $\tilde{T}$.
\end{Not}

\begin{Par}[Special State Functions]
The meaning of the special state functions will become clearer in \thref{def:canonical-state-functions} when we define their canonical form. As a brief introduction, a network function obviously determines the state of an entire network. On the other hand, a step of a neural network resembles a single step in the development of the network towards the final state of a network function. An important property of a sensible network function is \emph{balancedness} which we will define below.
\end{Par}

\begin{Def}[Balancedness]\label{def:balancedness}
A state $s\in S[\mathfrak{N}]$ is \emph{balanced} if and only if 
\begin{equation}\label{eq:balancedness}
s^{:o}=f(s^{:i}),
\quad
s^{:i}=g_{\theta,\lambda}(s^{:o}).
\end{equation}
We define
\[
B[\mathfrak{N}]:=\left\{s\in S[\mathfrak{N}]:s\text{ balanced}\right\}.
\]
A network function $T$ is balanced if and only if for all 
\[
T\left(S^{\mathcal{I}:i}_{\const,\var}[\mathfrak{N}]\right)\subseteq B[\mathfrak{N}].
\]
A state function $T$ is balanced if and only if there is a balanced network function $\tilde{T}$ such that
\[
\tilde{T}^{\mathcal{I}}=T.
\]
\end{Def}

\begin{Def}[Canonical State Functions]\label{def:canonical-state-functions}
Consider a network $\mathfrak{N}=(\mathcal{V},\theta)$. Consider a topology $\mathcal{T}$ on $A$.
\begin{enumerate}[a)]
\item
The \emph{canonical step} $T$ for a state $s\in S^{\mathcal{V}:i}_{\const,\var}[\mathfrak{N}]$ is given by
\begin{equation}\label{eq:canonical-step}
T_{\theta,\lambda}^{:o}(s):=T_{\theta,\lambda}^{:o}[\mathfrak{N}](s):=f\left(s^{:i}\right),
\quad
T_{\theta,\lambda}^{:i}(s):=T_{\theta,\lambda}^{:i}[\mathfrak{N}](s):=g_{\theta,\lambda}\left(T_{\theta,\lambda}^{:o}(s)\right).
\end{equation}
We omit the specification of the neural network $\mathfrak{N}$ if it is clear.
\item
The \emph{canonical network function} $N:=\mathfrak{N}$ for a state $i\in S^{\le\mathcal{V}:i}_{\const,\var}[\mathfrak{N}]$ where, for any $V\subseteq\mathcal{V}$, 
\[
S^{\le V}_{\const,\var}[\mathfrak{N}]:=\bigcup_{W\subseteq V}S^W_{\const,\var}[\mathfrak{N}],
\]
is given by
\begin{equation}
N_{\theta,\lambda}[\mathfrak{N}](i):=\lim_{n\to\infty}T_{\theta,\lambda}^n[\mathfrak{N}](\In(i))
\end{equation}
and is only defined for values where this limit exists in $(S[\mathfrak{N}],\mathcal{T}^{S[\mathfrak{N}]})$ where $\mathcal{T}^{S[\mathfrak{N}]}$ is the product topology on $S[\mathfrak{N}]$. We again omit the specification of the neural network if it is clear.
\end{enumerate}
\end{Def}

\begin{Par}[Convention: Real Neural Networks]
From now on, we will by convention restrict ourselves to real neural networks, i. e. $A=\mathbb{R}$ where $\mathcal{T}$ is the norm topology.
\end{Par}

\begin{Lem}[Fixed States]\label{lem:fixed-states}
Let $\text{FP}_T$ be the set of fixed points of $T$, i. e.
\[
\text{FP}_T:=\left\{s\in S[\mathfrak{N}]:T_{\theta,\lambda}(s)=T_{\theta,\lambda}\right\}.
\]
In this case,
\[
B[\mathfrak{N}]=\text{FP}_T\subseteq N_{\theta,\lambda}\left(S_{\const,\var}[\mathfrak{N}]\right)=N_{\theta,\lambda}\left(S^{\le\mathcal{V}}_{\const,\var}[\mathfrak{N}]\right)
\]
and, if $T$ is continuous, 
\[
\text{FP}_T=N\left(S_{\const,\var}[\mathfrak{N}]\right)
\]
\end{Lem}

\begin{Bew}
$B[\mathfrak{N}]=\text{FP}_T$: The fixed point equation corresponding to (\ref{eq:canonical-step}) is (\ref{eq:balancedness}).\\
$\text{FP}_T\subseteq N_{\theta,\lambda}\left(S_{\const,\var}[\mathfrak{N}]\right)$: If $s\in S[\mathfrak{N}]$ is a fixed point of $T$, then
\[
N_{\theta,\lambda}(s)=\lim_{n\to\infty}T_{\theta,\lambda}^n(s)=\lim_{n\to\infty}s=s.
\]
 $N_{\theta,\lambda}\left(S_{\const,\var}[\mathfrak{N}]\right)\subseteq N_{\theta,\lambda}\left(S^{\le\mathcal{V}}_{\const,\var}[\mathfrak{N}]\right)$ follows directly from $S_{\const,\var}[\mathfrak{N}]\subseteq S_{\const,\var}^{\le\mathcal{V}}[\mathfrak{N}]$.\\
On the other hand, for any $i\in S^{\le\mathcal{V}}_{\theta,\const}[\mathfrak{N}]$, $N_{\theta,\lambda}(i)=N_{\theta,\lambda}\left(\In(i)\right)$.\\
$\text{FP}_T\supseteq N\left(S_{\const,\var}[\mathfrak{N}]\right)$: Suppose that $S$ is continuous. Then, for any $i\in S_{\const,\var}[\mathfrak{N}]$:
\begin{align*}
N_{\theta,\lambda}(i)=\lim_{n\to\infty}T_{\theta,\lambda}^n(i)=T_{\theta,\lambda}\left(\lim_{n\to\infty}T^n_{\theta,\lambda}(i)\right)=T_{\theta,\lambda}\left(N_{\theta,\lambda}(i)\right).
\end{align*}
Therefore, $N_{\theta,\lambda}(i)$ is a fixed point of $T$.
\end{Bew}

\begin{Par}[Constant Values and Network Functions]
We often want the input to remain constant when applying a network function $T$. This is the reason why $T$ is defined for extended values.
\end{Par}

\begin{Def}[Implementations]
Consider a neural network $\mathfrak{N}:=(\mathcal{V},\theta)$, $\mathcal{I},\mathcal{O}\subseteq\mathcal{V}$, and a state function
\[
T_1:S^{\mathcal{I}}_{\const,\var}[\mathfrak{N}]\to S^{\mathcal{O}}_{\const,\var}[\mathfrak{N}].
\]
\begin{enumerate}[a)]
\item
If $\mathcal{O}'\supseteq\mathcal{O}$, 
\[
T_2:S^{\mathcal{O}'}_{\const,\var}[\mathfrak{N}]\to S^{\mathcal{O}'}_{\const,\var}[\mathfrak{N}]
\]
\emph{describes} $T_1$ if and only if the coordinates in $\mathcal{O}$ of the fixed points of $T_2$ are given by the image of $T_1$.
\item
A neural network $\tilde{\mathfrak{N}}:=(\tilde{\mathcal{V}},\tilde{\theta})$ with functional structure $(\tilde{\Lambda},\tilde{f},\tilde{g})$ such that
\[
\tilde{\mathcal{V}}\subseteq\mathcal{V},
\quad
\tilde{\theta}_{\mathcal{V}\to\mathcal{V}}=\theta,
\]
is called \emph{implementation} of $T_1$ if and only if its canonical step $T_{\tilde{\theta},\tilde{\lambda}}[\tilde{\mathfrak{N}}]$ describes $T_1$.
\end{enumerate}
\end{Def}

\subsection{States in Feed-Forward and Additive Neural Networks}\label{subsec:theory-states-add}

\begin{Par}
This subsection will clarify the rather general definitions of the previous paragraphs. More specifically, the canonical network function and step have a more concise form for feed-forward and additive networks. The following subsection \ref{subsec:theory-states-predicode} will make the use of these definitions apparent as they allow us to talk more precisely about implementations.
\end{Par}

\begin{Pro}
Consider a feed-forward network $\mathfrak{N}=(\mathcal{V},\theta)$ with functional structure $(\Lambda,f,g)$ of depth $L$.
If $i\in S^{\le\mathcal{V}:i}[\mathfrak{N}]$ is constant in the standardized layers $V^{(0)},\dotsc,V^{(l-1)}$, 
\begin{equation}
N_{\theta,\lambda}^{:i}(i)=\left(T_{\theta,\lambda}^{L-l+1}\right)^{:i}\left(\In(i)\right),
\quad
N_{\theta,\lambda}^{:o}(i)=f\left(N_{\theta,\lambda}^{:i}(i)\right).
\end{equation}
In particular, 
\begin{equation}
N_{\theta,\lambda}^{:i}=\left(T_{\theta,\lambda}^{L+1}\right)^{:i}\circ\In, 
\quad
N_{\theta,\lambda}^{:o}=f\left(N_{\theta,\lambda}^{:i}(i)\right).
\end{equation}
\end{Pro}

\begin{Bew}
We prove the proposition by proving the statement 
\[
\left(T_{\theta,\lambda}^{V^{(\le l)}:i}(i)=i^{V^{(\le l)}}\right)
\Rightarrow
\left(
\underbrace{\left(
N_{\theta,\lambda}^{:i}(i)=\left(T_{\theta,\lambda}^{L-l+1}\left(\In(i)\right)\right)^{:i}
\right)}_{(1)}
\wedge
\underbrace{\left( 
N_{\theta,\lambda}^{:o}(i)=f\left(N_{\theta,\lambda}^{:i}(i)\right)
\right)}_{(2)}
\right)
\]
for any $i\in S^{\le\mathcal{V}:i}[\mathfrak{N}]$ by induction on $L-l$. Note that the statement is slightly more general than the proposition.\\
If $L-l+1=0$, i. e. $l=L+1$, (1) follows directly and (2) is obvious from the definition of $N$.\\
Next, suppose that the proposition is proven for $L-l=L-(l+1)+1$ and consider some $i$ such that
\[
T_{\theta,\lambda}^{V^{(\le l)}}(i)=i^{V^{(\le l)}}.
\]
In this case, define
\[
s:=T_{\theta,\lambda}\left(\In(i)\right)
\]
We observe that
\[
s^{V^{(\le l)}}=i^{V^{(\le l)}},
\quad
s^{V^{(l+1)}}=T_{\theta,\lambda}^{V^{(l+1)}}\left(\In(i)\right)=T_{\theta,\lambda}^{V^{(l+1)}}\left(T_{\theta,\lambda}\left(\In(i)\right)\right)=T_{\theta,\lambda}^{V^{(l+1)}}(s),
\]
as the network is feed-forward and $T_{\theta,\lambda}^{V^{(l+1)}}$ therefore only depends on $i^{V^{(\le l)}}$. This implies
\begin{align*}
T_{\theta,\lambda}^{V^{(\le l+1)}}(s)=s^{V^{(\le l+1)}}
\end{align*}
and thus, by induction hypothesis, 
\[
N_{\theta,\lambda}^{:i}(i)=N_{\theta,\lambda}^{:i}(s)=\left(T_{\theta,\lambda}^{L-(l+1)+1}\left(\In(s)\right)\right)^{:i}.
\]
As $s\in S[\mathfrak{N}]$, $\In(s)=s$ and we therefore obtain
\[
N_{\theta,\lambda}^{:i}(i)=\left(T_{\theta,\lambda}^{L-l}\left(T_{\theta,\lambda}\left(\In(i)\right)\right)\right)^{:i}=\left(T_{\theta,\lambda}^{L-l+1}\left(\In(i)\right)\right)^{:i},
\]
which concludes the induction and thus proves the proposition.
\end{Bew}

\begin{Par}
In an additive neural network, the canonical step function is defined by
\begin{equation}
T_{\theta}=\theta^Tf
\end{equation}
and therefore, the canonical network function is defined by 
\begin{equation}
N_{\theta}=\lim_{n\to\infty}(\theta^T f)^n\circ\In.
\end{equation}
In particular, if $L$ is the depth of the network, the canonical network function of an additive feed-forward network is given by
\[
N_{\theta}^{:i}=(\theta^T f)^{L+2}\circ\In,
\quad
N_{\theta}^{:o}=f.
\]
\end{Par}

\subsection{States in Predictive Coding}\label{subsec:theory-states-predicode}

\begin{Par}[Predictive Coding II: The Free Energy Function]
Recall the generative model from \thref{def:generative-model}. As it is impossible to infer the global parameter $\omega\in\Omega$, we may consider the alternative functional structure $(\Lambda')$. However, we are now missing $g$ and therefore the canonical step and the canonical network function. Our goal is to infer the states $s\in S[\mathfrak{G}]$ from some input $i\in S^{\mathcal{I}}[\mathfrak{G}]$ and find a suitable implementation for this inference.\\
The inferred values should be as compatible as possible. We therefore posit a loss function
\begin{equation}
L_{\xi,\lambda'}:S[\mathfrak{G}]\to\mathbb{R}
\end{equation}
which we want to minimize. The knowledge we possess about how the states of the individual PUs should relate to each other is given by $\xi$ and $\lambda'$ which are treated as fixed parameters.\\
The input values $i\in S^{\mathcal{I}}[\mathfrak{G}]$ are already known and should therefore remain constant. The ideal network function would therefore return
\begin{equation}
\arg\min_{s\in S^{\mathcal{G}\setminus\mathcal{I}}[\mathfrak{G}]}L_{\xi,\lambda'}(s)
\end{equation}
If we suppose a continuous random variable with density $p_{\xi,\lambda'}$ a sensible choice of $L$ is the negative log-likelihood:
\begin{equation}
L_{\xi,\lambda'}(s):=-\ln p_{\xi,\lambda'}(s)
\end{equation}
The negative log-likelihood represents our surprise about the data. It is also called \emph{free energy}.
\end{Par}

\begin{Pro}[Free Energy in the Normal Generative Model]\label{pro:free-energy-normal}
If $L_{\xi,\lambda'}$ is the free energy function in the normal generative model and
\begin{equation}
F_{\xi,\lambda'}(s)=
\frac12\left(\ln\left(\det\Sigma_{\lambda'}\right)+\epsilon^{T}\Sigma_{\lambda'}^{-1}\epsilon\right)
\end{equation}
where
\[
\epsilon=s^{:i}-\mu_{\xi,\lambda'}(s^{:o}), 
\]
the folllowing identities hold\footnote{
The right part evidently only holds if $L$ is differentiable.}
\begin{align}
\arg\min_{s\in S[\mathfrak{G}]}F_{\xi,\lambda'}(s)&=\arg\min_{s\in\mathbb{R}^{\mathcal{G}\setminus\mathcal{I}}}L_{\xi,\lambda'}(s),&
\frac{\partial}{\partial s}F_{\xi,\lambda'}(s)&=\frac{\partial}{\partial s}L_{\xi,\lambda'}(s),\\
\arg\min_{\xi}F_{\xi,\lambda'}(s)&=\arg\min_{\xi}L_{\xi,\lambda'}(s),&
\frac{\partial}{\partial \xi}F_{\xi,\lambda'}(s)&=\frac{\partial}{\partial \xi}L_{\xi,\lambda'}(s),\\
\arg\min_{\lambda'}F_{\xi,\lambda'}(s)&=\arg\min_{\lambda'}L_{\xi,\lambda'}(s),&
\frac{\partial}{\partial \lambda'}F_{\xi,\lambda'}(s)&=\frac{\partial}{\partial \lambda'}L_{\xi,\lambda'}(s).
\end{align}
In other words: Minimizing $L$ corresponds to minimizing $F$.
\end{Pro}

\begin{Bew}
As we consider the normal generative model, $\epsilon$ is normally distributed:
\[
\epsilon\sim\mathcal{N}(0,\Sigma_{\lambda'})
\]
Therefore, 
\[
p_{\xi,\lambda'}(s)=\frac{1}{\sqrt{(2\pi)^{|\mathcal{G}|}\det\Sigma_{\lambda'}}}\exp\left(-\frac12\epsilon^T\Sigma^{-1}\epsilon\right)
\]
Therefore, 
\[
L_{\xi,\lambda'}(s)=-\ln p_{\xi,\lambda'}(s)=\ln\left(\sqrt{(2\pi)^{|\mathcal{G}|}}\right)+\ln\left(\sqrt{\det\Sigma_{\lambda'}}\right)+\frac12\epsilon^T\Sigma_{\lambda'}^{-1}\epsilon
\]
As the summand $\ln\left(\sqrt{(2\pi)^{|\mathcal{G}|}}\right)$ is constant in $s,\xi$ and $\Sigma$, omitting it does not change the minimum or the derivative in these variables. Observing
\[
\ln\left(\sqrt{\det\Sigma}\right)=\frac12\left(\ln\det\Sigma\right)
\]
yields $F$.
\end{Bew}

\begin{Par}[Predictive Coding III: Inference]
Generally, there is not a good method to find the global minimum of $L$. Instead, we propose optimization by gradient descent, such that a single step is given by $\varphi_{\xi,\lambda'|\alpha}\equiv\left(\phi_{\xi,\lambda'|\alpha}\right)_{\const,\var}$ where
\begin{equation}\label{eq:phi}
\phi_{\xi,\lambda'|\alpha}:S[\mathfrak{G}]\to S[\mathfrak{G}],
\quad
\phi_{\xi,\lambda'|\alpha}^{:i}(s):=s-\alpha\left(\frac{\partial L_{\xi,\lambda'}(s)}{\partial s}\right)^T,
\quad
\phi_{\xi,\lambda'|\alpha}^{:o}(s):=f\left(\phi_{\xi,\lambda'|\alpha}^{:i}(s)\right).
\end{equation}
The network function is therefore given by iterative gradient descent starting at $a$:
\begin{equation}
\Phi_{\xi,\lambda'|\alpha,a}:=\lim_{n\to\infty}\phi_{\xi,\lambda'|\alpha}^n\circ\In_a.
\end{equation}
Choosing $\alpha$ and $a$ as fixed parameters, we call an implementation of $\Phi_{\xi,\lambda'}$ a \emph{predictive coding network}.
\end{Par}

\begin{Not}[Derivatives of and in Vectors and Matrices]\label{matrix-derivation-not}
Given functions
\[
f:\mathbb{R}^m\to\mathbb{R}^n,
\quad
M:\mathbb{R}\to\mathbb{R}^{m\times n},
\quad
y:\mathbb{R}^{m\times n}\to\mathbb{R},
\]
we notate the derivatives by
\begin{align*}
\frac{\partial f}{\partial x}=\left(\frac{\partial f_i}{\partial x_j}\right)_{\substack{i=1,\dotsc,m\\j=1,\dotsc,n}},
\quad
\frac{\partial M}{\partial x}=\left(\frac{\partial M_{ij}}{\partial x}\right)_{\substack{i=1,\dotsc,m\\j=1,\dotsc,n}},
\quad
\frac{\partial y}{\partial A}=\left(\frac{\partial y}{\partial A_{ji}}\right)_{\substack{i=1,\dotsc,n\\j=1,\dotsc,m}}.
\end{align*}
Appendix \ref{ch:matrix-derivation}
\end{Not}

\begin{Lem}[Gradient Descent in a Normal Generative Model]\label{lem:grad-desc-normal}
In a normal generative model, if $\mu$ is differentiable, the step is given by
\begin{equation}\label{eq:normal-gradient-descent}
\phi_{\xi,\Sigma|\alpha}(s):=s-
\alpha\Sigma^{-1}\epsilon
+
\alpha\left(\frac{\partial}{\partial s}\mu_{\xi}(s)\right)^T\Sigma^{-1}\epsilon
\in 
S[\mathfrak{G}].
\end{equation}
\end{Lem}

\begin{Bew}
According to the chain rule,
\begin{align*}
\frac{\partial F}{\partial s}=\frac{\partial F}{\partial\epsilon}\frac{\partial\epsilon}{\partial s},
\end{align*}
\thref{derivation-quadratic-form} yields
\[\frac{\partial F}{\partial\epsilon}=\frac{\partial\frac12\epsilon^T\Sigma^{-1}\epsilon}{\partial\epsilon}=\epsilon^T\Sigma^{-1}.\]
On the other hand, 
\begin{align*}
\frac{\partial\epsilon}{\partial s}=
\frac{\partial}{\partial s}(s-\mu_{\xi}(s))=
I_{\mathcal{G}\to\mathcal{G}}-\frac{\partial\mu_{\xi}(s)}{\partial s}.
\end{align*}
Putting both together yields
\[
\frac{\partial F}{\partial s}=
\Sigma^{-1}\epsilon
-
\left(\left(\frac{\partial}{\partial s}\mu_{\xi}(s)\right)^T\Sigma^{-1}\epsilon\right)^T,
\]
which, if filled into (\ref{eq:phi}), yields (\ref{eq:normal-gradient-descent}).
\end{Bew}

\begin{Par}
From now on, we will specialize in additive networks as many of the equations have a rather unintuitive implementation whose notation seems forced.
\end{Par}

\begin{Cor}[Gradient Descent in an Additive Normal Generative Model]\label{grad-desc-add}
In a normal additive generative model, the step is given by
\begin{equation}\label{eq:grad-desc-add}
\phi_{\xi,\Sigma|\alpha}(s):=s-\alpha\Sigma^{-1}\epsilon
+\alpha
f'(s)\xi\Sigma^{-1}\epsilon,
\end{equation}
where
\begin{equation}
f':=\frac{\partial}{\partial s}f.
\end{equation}
\end{Cor}

\begin{Bew}
By the definition of additive networks, 
\[
\mu_{\xi}(s)=\xi^Tf(s),
\]
which implies
\begin{align*}
\frac{\partial}{\partial s}\mu_{\xi}(s)=\xi^T\frac{\partial}{\partial s}f(s)=\xi^Tf'(s).
\end{align*}
The corollary now follows directly from \thref{lem:grad-desc-normal} and the fact that $f'(s)$ is diagonal and therefore symmetric.
\end{Bew}

\begin{The}[Additive Normal Predictive Coding Network]\label{the:add-predicode}
Gradient descent in a normal additive generative model $\mathfrak{G}=(\mathcal{G},\xi)$ with a functional structure $(\Lambda'\times\Omega,\tilde{f})$ such that $\tilde{f}$ is continuous can be implemented by a neural network $\mathfrak{P}=(\mathcal{P},\theta)$ where
\begin{equation}\label{eq:add-predicode}
\begin{aligned}
\mathcal{P}&=\mathcal{G}\cup\varepsilon,
\quad
\varepsilon=\left\{\varepsilon^v|v\in\mathcal{G}\right\}\\
\theta_{v\to w}&:=
\begin{cases}
(\delta_{v,w},\const) &\text{ if }v,w\in\mathcal{G},\\
(\delta_{v,u},\const)-\xi_{v\to u} &\text{ if }v\in\mathcal{G},w=\varepsilon^u\in\varepsilon,\\
\alpha\left(-(\delta_{u,w},\const)+\xi^T_{u\to w} \right) &\text{ if }v=\varepsilon^u\in\varepsilon,w\in\mathcal{G},\\
(\delta_{u,x},\const)-\left(\Sigma_{\lambda'}^T\right)_{u\to x} &\text{ if }v=\varepsilon^u,w=\varepsilon^x\in\varepsilon.
\end{cases}
\end{aligned}
\end{equation}
Identifying the entries in $\mathcal{G}$ and $\varepsilon$, we may also use a shorthand matrix notation:
\begin{align*}
\theta=
\begin{pmatrix}
\theta_{\mathcal{G}\to\mathcal{G}} & \theta_{\mathcal{G}\to\varepsilon}\\
\theta_{\varepsilon\to\mathcal{G}} & \theta_{\varepsilon\to\varepsilon}
\end{pmatrix}
=
\begin{pmatrix}
I_{\const} & I_{\const}-\xi\\
-\alpha I_{\const}+\alpha\xi^T & (I-\Sigma)^T
\end{pmatrix}.
\end{align*}
The functional structure of $\mathfrak{P}$ is given by $(f,g)$, where
\begin{equation}
\begin{aligned}
f:S^{:i}[\mathfrak{P}]&\to S^{:o}[\mathfrak{P}]
&
s^{v:i}&\mapsto s^{v:o}:=f^v(s^{v:i}):=\begin{cases}
\tilde{f}^v(s^{v:i})&\text{ if }v\in\mathcal{G},\\
s^{v:i}&\text{ if }v\in\varepsilon,
\end{cases}\\
g_{\theta}:S[\mathfrak{P}]&\to S^{:i}[\mathfrak{P}]
&
s&\mapsto s^{v:i}:=g^v_{\theta}(s):=\begin{cases}
s^{v:i}+\theta_{\varepsilon^v\to v}s^{\varepsilon^v:o}+\left(f^v\right)'(s^{v:i})\theta_{\varepsilon\setminus\{\varepsilon^v\}\to v}^Ts^{\varepsilon\setminus\{\varepsilon^v\}:o}&\text{ if }v\in\mathcal{G},\\
\theta_{\mathcal{P}\setminus\{u\}\to v}^Ts^{\mathcal{P}\setminus\{u\}:o}+\theta_{u\to v}s^{v:i}&\text{ if }v=\varepsilon^u\in\varepsilon.
\end{cases}
\end{aligned}
\end{equation}
We refer to $\varepsilon$ as the \emph{error units} and $\mathcal{G}$ as the generative units.
\end{The}

\begin{Par}
Evidently, the additive normal predictive coding network is not a general additive network. However, the only reason for this is that the input from the other error units is weighted by the derivative of the generative unit function. (Note that the neural network is still well-defined as there is an auto-connection $v\to v$.) It is possible to implement the gradient descent in a general additive network by adding derivative nodes and setting $\ln$ and $\exp$ as suitable unit functions. As long as there is no apparent use for such a specification, we will refrain from such a definition.
\end{Par}

\begin{Bew}[of \thref{the:add-predicode}]
As mentioned, $\mathfrak{P}$ is indeed a neural network. The additive dependency structure makes that immediately visible if we note that $\left(f^v\right)'(s^{v:i})$ is known due to the connection $v\to v$.\\
In order to show that $\mathfrak{P}$ implements the gradient function, we have to characterize the fixed points of the canonical step function given by the equation
\[
g_{\theta}(s)=s^{:i},
\quad
f(s^{:i})=s^{:o}.
\]
We first consider the error units $\varepsilon^v\in\varepsilon$:
\[
s^{\varepsilon^v:i}=g_{\theta}(s^{\varepsilon^v})=\theta_{\mathcal{P}\setminus\{v\}\to \varepsilon^v}^Ts^{\mathcal{P}\setminus\{v\}:o}+\theta_{v\to \varepsilon^v}s^{\varepsilon^v:i}\overset{(*)}{=}-\xi^Ts^{\mathcal{G}\setminus\{v\}:o}+s^{v:i}+(I-\Sigma)s^{\varepsilon:o}=\epsilon^v+(I-\Sigma)s^{\varepsilon:o},
\]
where $(*)$ is implied by the fact that $\xi$ is an upper triagonal matrix $\xi^T$ therefore has an empty diagonal. As $s^{\varepsilon:i}=s^{\varepsilon:o}=:s^{\varepsilon}$, we may conclude
\[
s^{\varepsilon}=\varepsilon+(I-\Sigma)s^{\varepsilon}=\Sigma^{-1}\epsilon
\]
Considering the generative units $v\in\mathcal{G}$ and using this equation, we can now conclude
\[
s^{v:i}=s^{v:i}+\theta_{\varepsilon^v\to v}s^{\varepsilon^v:o}+\left(f^v\right)'(s^{v:i})\theta_{\varepsilon\setminus\{\varepsilon^v\}\to v}^Ts^{\varepsilon\setminus\{\varepsilon^v\}:o}=s^{v:i}-\alpha s^{\varepsilon^v}+\alpha\left(f^v\right)'(s^{v:i})\xi_{v\to} s^{\varepsilon}
\]
where we have again used that the diagonal of $\xi$ is empty. Vectorizing the equation, we observe
\[
s^{:i}=s^{:i}-\alpha \Sigma^{-1}\epsilon+\alpha\left(\frac{\partial}{\partial s^{:i}}f(s^{:i})\right)\xi\Sigma^{-1}\epsilon
\]
This corresponds to the step function $\phi_{\xi,\Sigma|\alpha}$. As $\tilde{f}$ is continuous, $\phi$ is continuous. As $\mathbb{R}$ is closed, the value $\Phi$ converges to are the fixed point of $\phi$ which are the $\mathcal{G}$-coordinates of the fixed points of the canonical step functions. If we are given an element of the image of $\Phi$, we may simply set $s^{\varepsilon}:=\Sigma^{-1}\epsilon$ which yields a fixed point of the canonical step function $T[\mathfrak{P}]$ of $\mathfrak{P}$. We have thus proven the theorem.
\end{Bew}

\begin{Lem}\label{lem:predicode-predistate}
Consider an additive normal predictive coding network $\mathfrak{P}=(\mathcal{G}\cup\varepsilon,\theta)$ with the functional structure $(f,g)$. If the input is $i\in S^{V^{(0)}}_{\const,\var}[\mathfrak{P}]$ and $N(i)$ exists, the state of its canonical network function $s:=N^{\mathcal{G}}[\mathfrak{P}](i)$ may be defined recursively in the standardized layerization $G^{(0)}\prec G^{(1)}\prec\dotsb\prec G^{(L)}$ by
\begin{equation}\label{eq:predicode-predistate}
\begin{aligned}
s^{v:i}&:=
\begin{cases}
\Sigma_{v\to G^{(0)}}s^{G^{(0)}:i}&\text{ if }v\in G^{(0)}\wedge\St(i^v)=\var,\\
i^v&\text{ if }v\in G^{(0)}\wedge\St(i^v)=\const,\\
\xi^Ts^{:o}+\Sigma_{v\to G^{(0)}}s^{G^{(0)}:i}&\text{ if }v\notin G^{(0)},
\end{cases}
\\
s^{v:o}&:=f(s^{v:i}).
\end{aligned}
\end{equation}
In particular, if
\[
\underset{v\in G^{(0)}}{\forall}\underset{w\in \mathcal{G}\setminus G^{(0)}}{\forall}\Sigma_{v\to w}=0,
\]
i. e. if the input nodes are independent from all other nodes and $i$ is constant, 
\begin{equation}
s^{:i}:=\left(\xi^Tf\right)^{L}\left(\In(i)\right),
\quad
s^{:o}:=f(s^{:i}).
\end{equation}
\end{Lem}

\begin{Bew}
Define $\tilde{s}:=N(i)$. Evidently, $\tilde{s}^{\mathcal{G}}=s$ and $\tilde{s}\in B[\mathfrak{P}]$. Consider $v\in G^{(0)}$. If $\St(i^v)=\const$, $\tilde{s}^{v:i}=i^v$ as $N$ is an active extension. If $\St(i^v)=\var$ as well as if $v\in\mathcal{G}\setminus G^{(0)}$, then the fixed point equation implies
\[
s^{v:i}=S^{v:i}(\tilde{s})=s^{v:i}-\alpha \tilde{s}^{\varepsilon^v}+\alpha\left(f^v\right)'(s^{v:i})\xi_{v\to} \tilde{s}^{\varepsilon}.
\]
This implies
\[
\tilde{s}^{\varepsilon^v}=\left(f^v\right)'(s^{v:i})\xi_{v\to}\tilde{s}^{\varepsilon}.
\]
We now show by induction on $\leftarrow$ that if $\St(i^v)=\var$ or $v\in\mathcal{G}\setminus G^{(0)}$,
\[
s^{\varepsilon^v}=0.
\]
Because $\xi$ is feed-forward and thus an upper triangular matrix where
\[
\xi_{v\to G^{(L)}}=0,
\]
we obtain, for all $v\in G^{(L)}$,
\[
\tilde{s}^{\varepsilon^v}=0,
\]
which proves the base case.
On the other hand, if
\[
\tilde{s}^{\varepsilon^{v\to}}=0,
\]
we can conclude that
\[
\tilde{s}^{\varepsilon^{v}}=\left(f^v\right)'(s^{v:i})\xi_{v\to(v\to)}\tilde{s}^{\varepsilon^{v\to}}=0,
\]
which proves the induction step.\\
As $\tilde{s}$ is balanced, we can also infer that
\[
\tilde{s}^{\varepsilon}=\Sigma^{-1}(s^{:i}-\xi^Ts^{:o})
\]
and thus
\[
s^{:i}=\Sigma\tilde{s}^{\varepsilon}+\xi^Ts^{:o}.
\]
As $\tilde{s}^{\varepsilon^{\mathcal{G}\setminus G^{(0)}}}=0$ and $\xi^{T}_{G^{(0)}\to}=0$, this implies (\ref{eq:predicode-predistate}).
\end{Bew}

\begin{Par}
\thref{lem:predicode-predistate} demonstrates that a common use of the predictive coding network where $V^{(0)}$ is the input layer and the input layer does not possess any correlation structure with the succeeding nodes\footnote{
Note that this is not a restriction as such a structure can be covered by a suitable $\xi$.} can be characterized by a standard additive neural network $(\mathcal{G},\xi)$ with functional structure $f$.
\end{Par}

\section{Learning in Neural Networks}\label{sec:theory-learning}

\begin{Par}
A crucial property of neural networks is \emph{learning}, their ability to change their parameters according to training data. Learning is introduced in subsection \ref{subsec:theory-learning-intro}. The backpropagation algorithm, a popular learning algorithm in feed-forward neural networks is introduced in \ref{subsec:theory-learning-backprop}. Subsection \ref{subsec:theory-learning-predicode} discusses biologically plausible learning and introduces the learning algorithm of a normal additive predictive coding network. The final subsection \ref{subsec:theory-learning-comp} provides some remarks on the relation between backpropagation and the predictive coding algorithm.
\end{Par}

\subsection{Learning Functions}\label{subsec:theory-learning-intro}

\begin{Def}[Learning Structure]
Consider a neural network $\mathfrak{N}=(\mathcal{V},\theta)$ with functional structure $(f,g,\Lambda)$. Let $\mathcal{I},\mathcal{O}\subseteq\mathcal{V}\times\{i,o\}$ be disjunct. The PUs $\mathcal{I}$ correspond to the input and $\mathcal{O}$ corresponds to the output that shall be approximated. The function
\[
\tilde{\mathcal{L}}:\left(S^{\mathcal{I}}[\mathfrak{N}]\times S^{\mathcal{O}}[\mathfrak{N}]\right)^{<\mathbb{N}}\times A^{\mathcal{V}\times\mathcal{V}}_{\const,\var}\times\Lambda\to A^{\mathcal{V}\times\mathcal{V}}_{\const,\var}\times\Lambda
\]
defines by its active extension in $A^{\mathcal{V}\times\mathcal{V}}_{\const,\var}$ a \emph{learning function}
\[
\mathcal{L}:=\tilde{\mathcal{L}}_{\const,\var:A^{\mathcal{V}\times\mathcal{V}}_{\const,\var}},
\quad
(\left((i_1,o_1),\dotsc,(i_n,o_n)\right),\theta,\lambda)\mapsto \mathcal{L}^{i;o}(\theta,\lambda)=(C^{i;o}(\theta,\lambda),G^{i;o}(\theta,\lambda)),
\]
where
\[
C^{i;o}(\theta,\lambda)\in A^{\mathcal{V}\times\mathcal{V}}
\]
returns the actualized connectionist pattern and
\[
G^{i;o}(\theta,\lambda)\in\Lambda
\]
returns the actualized global parameters.
\end{Def}

\begin{Par}\label{training-extension}
We will mostly consider a learning function for a single training example. In many cases, an extension towards several training examples can be given by
\begin{equation}
\mathcal{L}^{i_1,\dotsc,i_n;o_1,\dotsc,o_n}_{\theta,\lambda}:=\mathcal{L}^{i_n;o_n}\left(\mathcal{L}^{i_{n-1};o_{n-1}}\left(\dotsb \mathcal{L}^{i_1;o_1}(\theta,\lambda)\dotsb\right)\right)
\end{equation}
\end{Par}

\begin{Par}[Learning by Gradient Descent]
An important class of learning functions is given by gradient descent functions which are induced by a loss function. Loss functions are defined for some $\mathcal{O}\subseteq\mathcal{V}\times\{i,o\}$ by
\begin{equation}
L:\left(S^{\mathcal{O}}[\mathfrak{N}]\times S^{\mathcal{O}}[\mathfrak{N}]\right)^{<\mathbb{N}}\to\mathbb{R},
\quad
\left((\hat{o}_1,o_1),\dotsc,(\hat{o}_n,o_n)\right)\mapsto L\left(\hat{o},o\right).
\end{equation}
The gradient descent is given by
\[
\tilde{\mathcal{L}}_{\alpha}^{i;o}(\theta,\lambda):=\left(\theta-\alpha\left(\frac{\partial L\left(N_{\theta,\lambda}(i);o\right)}{\partial\theta}\right)^T,
\lambda-\alpha\left(\frac{\partial L\left(N_{\theta,\lambda}(i);o\right)}{\partial\lambda}\right)^T\right)
\]
such that the gradient descent learning function is defined by its active expansion in $\mathbb{R}^{\mathcal{V}\times\mathcal{V}}$:
\begin{equation}
\mathcal{L}_{\alpha}:=\left(\tilde{\mathcal{L}}_{\alpha}\right)_{\const,\var:\mathbb{R}^{\mathcal{V}\times\mathcal{V}}}
\quad
\alpha>0
\end{equation}
$\mathcal{L}_{\alpha}$ is \emph{induced} by $L$.
\end{Par}

\subsection{Backpropagation}\label{subsec:theory-learning-backprop}

\begin{Par}
Many feed-forward neural networks are purely concerned with prediction: the lowest layer $V^{(0)}$ represents the input and the highest layer $V^{(L)}$ represents the output that is predicted. In this case, we notate
\begin{equation}
O:=N^{V^{(L)}:o}.
\end{equation}
If $i$ is the input and $o$ is the actual result, we now try to optimize some loss function
\[
L\left(O(i),o\right).
\]
A popular choice for $L$ is for instance the squared error. The backpropagation algorithm can compute the gradient descent learning function and is especially suited for additive feed-forward networks. We therefore present backpropagation in this framework below.
\end{Par}

\begin{The}[Backpropagation]\label{the:backpropagation}
Consider an additive feed-forward network $\mathfrak{N}=(\mathcal{V},\theta)$ with the functional structure $f$ and the loss function 
\begin{equation}
L:\mathbb{R}^{V^{(L)}}\times \mathbb{R}^{V^{(L)}}\times\mathbb{R}_{\const,\var}^{\mathcal{V}\times\mathcal{V}}\to\mathbb{R}.
\end{equation}
Furthermore, suppose that there are no PUs $u,v,w$ such that $u\to w$ and $u\to v\to^*w$.\\
The gradient can then be computed by
\begin{equation}\label{eq:backprop-1}
\frac{\partial L(O_{\theta}(i),o)}{\partial\theta_{v\to w}}=\delta^wN^{v:o}(i)
\end{equation}
where
\[
\delta:=\frac{\partial L}{\partial N^{v:i}_{\theta}(i)}
\]
is recursively defined by
\begin{equation}\label{eq:backprop-2}
\begin{aligned}
v\in V^{(L)}&\Rightarrow\delta^v=\left(f^v\right)'\left(N^{v:i}_{\theta}(i)\right)\frac{\partial L_{\theta}(O_{\theta}(i),o)}{\partial O_{\theta}^v(i)}\\
v\notin V^{(L)}&\Rightarrow\delta^v=\left(f^v\right)'\left(N^{v:i}_{\theta}(i)\right)\delta^{v\to}\theta_{v\to(\to v)}.
\end{aligned}
\end{equation}
We define the learning function $\mathcal{B}_{\alpha}$ that is induced by $L$.
\end{The}

\begin{Bew}
By
\[
\frac{\partial L(O_{\theta}(i),o)}{\partial\theta_{v\to w}}=\frac{\partial L(O_{\theta}(i),o)}{\partial N_{\theta}^{w:i}(i)}\frac{\partial N_{\theta}^{w:i}(i)}{\partial\theta_{v\to w}}
\]
we only have to note that
\[
\frac{\partial}{\partial\theta_{v\to w}}N^{w:i}(i)=\frac{\partial}{\partial\theta_{v\to w}}\theta^T_{\to w}N^{:o}(i)=N^{v:o}(i)
\]
to get (\ref{eq:backprop-1}).\\
We therefore prove that (\ref{eq:backprop-2}) defines $\delta$ by induction. If $v\in V^{(L)}$, i. e. $v\to=\emptyset$,
\[
N^{v:i}_{\theta}(i)=O^v_{\theta}(i)
\]
and (\ref{eq:backprop-2}) is therefore obvious. On the other hand, if we already know
\[
\delta^{v\to}=\frac{\partial L(O_{\theta}(i),i)}{\partial N^{v\to:i}_{\theta}(i)},
\]
we may write
\[
\delta^v=\delta^{v\to}\frac{\partial N^{v\to:i}_{\theta}(i)}{\partial N^{v:o}_{\theta}(i)}\frac{\partial N^{v:o}_{\theta}(i)}{\partial N^{v:i}_{\theta}(i)}.
\]
As
\begin{align*}
N^w_{\theta}(i)=\theta_{\to w}^TN_{\theta}(i)
\end{align*}
and for any $u\in\to w$, $N^u$ does not depend on $N^v$, 
\[
\frac{\partial N^{v\to:i}_{\theta}(i)}{\partial N^{v:o}_{\theta}(i)}=\theta_{v\to(\to w)},
\]
which, combined with the upper equation and the observation that
\[
\frac{\partial N^{v:o}_{\theta}(i)}{\partial N^{v:i}_{\theta}(i)}=\left(f^v\right)'(N^{v:i}_{\theta}(i)),
\]
yields (\ref{eq:backprop-2}).
\end{Bew}

\begin{Par}
The method the proof of the theorem sketches works for any feed-forward network. However, in the case of non-additive networks, it is not very enlightening.
\end{Par}

%\begin{The}[Backpropagation]\label{the:backpropagation}
%Consider an additive feed-forward network $\mathcal{N}=(\mathcal{V},\theta)$ with the functional structure $f$ and the loss function
%\begin{equation}
%L:\mathbb{R}^{V^{(0)}\cup V^{(L)}}\times\mathbb{R}_{\const,\var}^{\mathcal{V}\times\mathcal{V}}\to\mathbb{R}
%\end{equation}
%The gradient can then be calculated recursively by
%\begin{equation}
%\begin{aligned}
%\tilde{\delta}&:=\frac{\partial L_{\theta}(s^{V^{(0)}},s^{V^{(L)}})}{\partial N_{\theta}(s^{V^{(L)}})}\\
%\delta&:=\frac{\partial f\left(N_{\theta}(s)\right)}{\partial N_{\theta}(s)}\theta^L\In(\tilde{\delta})\\
%\frac{\partial L}{\partial\theta}&:=\delta f\left(N_{\theta}(s)\right)^T
%\end{aligned}
%\end{equation}
%\end{The}
%
%\begin{Bew}
%We observe that the definition of $\delta$, as $\theta$ is an upper triagonal matrix corresponds to the recursive definition
%\begin{align*}
%\hat{\delta}^{V^{(L)}}&:=\tilde{\delta}=\frac{\partial L}{\partial N^{V^{(L)}}}\\
%\underset{0\le l<L}{\forall}\delta^{V^{(l)}}&:=
%\end{align*}
%\end{Bew}

\subsection{Learning in Predictive Coding}\label{subsec:theory-learning-predicode}

\begin{Par}[Predictive Coding IV: Learning Perspectives]
In the previous paragraphs on predictive coding, we have constructed the generative model $\mathfrak{G}$, its parameters $\xi$ and $\lambda'$ (resp. $\Sigma_{\lambda'}$) and inference in $\mathfrak{G}$. Then came the implementation of inference: the predictive coding model $\mathfrak{P}$ and its parameters $\theta$. As $\theta$ depends on $\xi$ and $\Sigma_{\lambda'}$, we have two options with respect to parameter learning: we may either learn $\xi$ and $\Sigma_{\lambda'}$ and thereby $\theta$ or we may learn $\theta$.\\
While every learning rule on $\xi$ and $\Sigma_{\lambda'}$ can be transformed into a learning rule on $\theta$ this transformation only preserves the Hebbian property if $\theta$ is still structured as in (\ref{eq:add-predicode}) (see chapter \ref{ch:further-work}). On the other hand, not every learning rule on $\theta$ can be transformed into a learning rule on $\xi$ and $\Sigma$ but if there exists such a transformation, it preserves the Hebbian property.\\
For now, we will consider learning $\xi$ and $\lambda'$. The approach we take is based on the expectation maximization algorithm \cite{Dempster1977} as implemented in \cite{Friston2005} and \cite{Whittington2017}: we want to minimize the free energy function in $\xi$ and $\lambda'$ for a balanced state $s\in B[\mathfrak{P}]$. We are therefore looking for the learning function $\mathcal{L}_{\alpha}[\mathfrak{P}]$ induced by $L$. The difference to backpropagation and most other learning algorithms is that we do not compare the real outcome $o$ to the outcome according to $N[\mathfrak{P}](i)$. Rather, we obtain a balanced state for input and output together. We therefore often consider $i\in S^{\mathcal{V}^{(0,L)}:i}[\mathfrak{P}]$ (although any kind of input may be used) and then set
\[
s:=N_{\xi,\lambda'}[\mathfrak{P}](i)
\]
Now, we try to reduce the free energy function $L_{\theta(\xi,\lambda')}(s)$\footnote{
I have written $\theta(\xi,\lambda')$ to make the dependence clearer.} in $\xi$ and $\lambda'$. We therefore do not try to change the parameters such that $N[\mathfrak{P}]_{\theta(\xi,\lambda')}(i)$ produces more compatible results. Rather, we obtain these estimates of our values and then update $\xi$ and $\lambda'$ to make them more compatible.\\
In order to do that, we need the gradient of $L_{\theta(\xi,\lambda')}(s)$ for balanced states in $\xi$ and $\lambda'$ which the following two theorems provide for additive normal generative models.
\end{Par}

\begin{The}[Learning $\xi$]\label{learning-xi}
Consider an additive normal generative model $\mathfrak{G}=(\mathcal{G},\xi)$ with functional structure $(\Lambda,f)$. If $L_{\xi,\lambda'}$ is the free energy function in $\mathfrak{G}$ and $s$ is balanced, we obtain
\[
\frac{\partial L_{\xi,\lambda'}(s)}{\partial\xi_{v\to w}}=-s^{\varepsilon^w:i}s^{v:o}
\]
\end{The}

\begin{Bew}
As we have observed in the proof of \thref{the:add-predicode}, 
\begin{align*}
s^{\varepsilon:i}=\Sigma^{-1}\epsilon.
\end{align*}
According to \thref{pro:free-energy-normal},
\[
\frac{\partial L_{\xi,\lambda'}(s)}{\partial\xi_{v\to w}}=\frac{\partial F_{\xi,\lambda'}(s)}{\partial\xi_{v\to w}}
\]
\thref{derivation-quadratic-form} implies that
\[
\frac{\partial F_{\xi,\lambda'}(s)}{\partial\epsilon}=(\Sigma^{-1}\epsilon)^T=\left(s^{\varepsilon}\right)^T
\]
On the other hand,
\[
\frac{\partial\epsilon^u}{\partial \xi_{v\to w}}=\begin{cases}
-s^{v:o}&\text{ if }u=w\\
0 &\text{ else}
\end{cases}
\]
which implies that
\[
\frac{\partial F_{\xi,\lambda'}(s)}{\partial\xi_{v\to w}}=
\frac{\partial F_{\xi,\lambda'}(s)}{\partial\epsilon}
\frac{\partial\epsilon}{\partial\xi_{v\to w}}=
-s^{\varepsilon^w:i}s^{v:o}
\]
and thus proves the theorem.
\end{Bew}

\begin{The}[Learning $\lambda'$]\label{learning-lambda}
Consider an additive normal generative model $\mathfrak{G}=(\mathcal{G},\xi)$ with functional structure $(\Lambda,f)$ where $\Lambda\subseteq\mathbb{R}^I$ is open (where $I$ is some index set). If $L_{\xi,\lambda'}$ is the free energy function in $\mathfrak{G}$ and $s$ is balanced, we obtain
\begin{equation}\label{eq:learning-lambda}
\frac{\partial L_{\xi,\lambda'}(s)}{\partial\lambda'_i}=
\frac12\tr\left(\left(\Sigma_{\lambda'}^{-1}-\left(s^{\varepsilon}\right)\left(s^{\varepsilon}\right)^T\right)\frac{\partial\Sigma_{\lambda'}}{\partial\lambda'_i}\right).
\end{equation}
In particular, if $\lambda'=\Sigma$, 
\begin{equation}\label{eq:learning-sigma}
\frac{\partial L_{\xi,\Sigma}(s)}{\partial\Sigma}=
\frac12\left(\Sigma^{-1}-\left(s^{\varepsilon}\right)\left(s^{\varepsilon}\right)^T\right)
.
\end{equation}
\end{The}

\begin{Bew}
We begin by evaluating
\begin{align*}
\frac{\partial}{\partial\lambda_i'}\ln\det\Sigma_{\lambda'}
\end{align*}
for some $i\in I$. As $\Sigma_{\lambda'}$ is invertible, \thref{jacobis-formula} implies
\begin{align*}
\frac{\partial}{\partial\lambda_i'}\ln\det\Sigma_{\lambda'}=\frac{1}{\det\Sigma_{\lambda'}}\frac{\partial}{\partial\lambda_i'}\det\Sigma_{\lambda'}=\tr\left(\Sigma_{\lambda'}^{-1}\frac{\partial}{\partial\lambda_i'}\Sigma_{\lambda'}\right).
\end{align*}
On the other hand,
\begin{align*}
&\frac{\partial}{\partial\lambda_i'}\epsilon^T\Sigma_{\lambda'}^{-1}\epsilon\overset{\thref{chain-rule}}{=}
\tr\left(\frac{\partial\epsilon^T\Sigma_{\lambda'}^{-1}\epsilon}{\Sigma^{-1}_{\lambda'}}\frac{\partial\Sigma^{-1}_{\lambda'}}{\partial\lambda_i'}\right)
\overset{\substack{\thref{derivation-matrix-quadratic-form}\\\thref{inverse-chain}}}{=}\\
&\tr\left(\epsilon\epsilon^T\cdot(-1)\cdot\Sigma_{\lambda'}^{-1}\frac{\partial\Sigma_{\lambda'}}{\partial \lambda_i'}\Sigma_{\lambda'}^{-1}\right)\overset{s\in B[\mathfrak{P}]}{=}
-\tr\left(\Sigma_{\lambda'} s^{\varepsilon}\left(s^{\varepsilon}\right)^T\frac{\partial\Sigma_{\lambda'}}{\partial\lambda_i'}\Sigma_{\lambda'}^{-1}\right)\overset{\tr(AB)=\tr(BA)}{=}\\
&-\tr\left(\Sigma_{\lambda'}^{-1}\Sigma_{\lambda'}s^{\varepsilon}\left(s^{\varepsilon}\right)^T\frac{\partial\Sigma_{\lambda'}}{\partial\lambda_i'}\right)=
-\tr\left(s^{\varepsilon}\left(s^{\varepsilon}\right)^T\frac{\partial\Sigma_{\lambda'}}{\partial\lambda_i'}\right).
\end{align*}
The fact that $\tr(A)+\tr(B)=\tr(A+B)$ and \thref{pro:free-energy-normal} proves (\ref{eq:learning-lambda}).\\
On the other hand, if $\lambda'=\Sigma_{\lambda'}$, we have
\begin{align*}
\frac{\partial}{\partial\Sigma_{w\to v}}L_{\xi,\Sigma}(s)=
\frac12\tr\left(
\underbrace{\left(\Sigma^{-1}-\left(s^{\varepsilon}\right)\left(s^{\varepsilon}\right)^T\right)}_
{=:A}
\underbrace{\frac{\partial\Sigma}{\partial\Sigma_{w\to v}}}_{=:B}\right),
\end{align*}
where
\begin{align*}
B_{u\to u'}=\begin{cases}
1&\text{ if }u=w\wedge u'=v,\\
0&\text{ else}.
\end{cases}
\end{align*}
Therefore,
\begin{align*}
\tr(AB)=\sum_{v\in\mathcal{G}}(AB)_{v\to v}=
\sum_{u'\in\mathcal{G}}\sum_{u\in\mathcal{G}}A_{u'\to u}B_{u\to u'}=
A_{v\to w},
\end{align*}
which implies
\begin{align*}
\frac{\partial}{\partial\Sigma_{w\to v}}L_{\xi,\Sigma}(s)=\frac12\left(\Sigma_{\lambda'}^{-1}-\left(s^{\varepsilon}\right)\left(s^{\varepsilon}\right)^T\right)_{v\to w}
\end{align*}
and thus proves (\ref{eq:learning-sigma}).
\end{Bew}

\begin{Par}[Biological Plausibility I: The Hebbian Postulate]
Why do we consider the more complicated predictive coding model instead of relying on a simple feed-forward model with backpropagation? On the one hand, a predictive coding model promises a more flexibel model specification (see chapter \ref{ch:further-work}), on the other hand, its definition is biologically motivated. In 1949, Donald Hebb, a Canadian psychologist, first presented his postulate on learning:
\begin{quote}
When an axon of cell A is near enough to excite a cell B and repeatedly or presistently takes part in firing it, some growth process or metabolic change takes place in one or both cells such that A's efficiency, as one of the cells firing B, is increased. \cite[][62]{Hebb1949}
\end{quote}
This means that the correlated activity of two synapses strengthens their synaptic connection. The mathematical formulation of Hebbian learning (see \cite{Gerstner2002}) includes several considerations. I have decided on the following two definitions. I will continue with some remarks on the relationship between backpropagation and predictive coding and their biological plausibility.
\end{Par}

\begin{Def}[Local and Hebbian Learning]
A learning function $\mathcal{L}$ is \emph{local} if and only if, $C^{i;o}_{v\to w}$ only depends on $N^v(i),N^w(i)$ and $\theta_{v\to w}$, i. e. for any $v,w\in\mathcal{V}$, we may write
\[
C^{i;o}_{v\to w}=\tilde{C}_{v\to w}\left(N^v(i),N^w(i),\theta_{v\to w}\right)
\]
for some function $\tilde{C}_{v\to w}$.\\
A local learning function is \emph{Hebbian} if and only if, for every $v,w\in\mathcal{V}$, $C^i_{v\to w}$ is monotonically increasing in $N^{v:o}(i)\cdot N^{w:i}(i)$.
\end{Def}

\subsection{Backpropagation and Predictive Coding}\label{subsec:theory-learning-comp}

\begin{Par}[Biological Plausibility II: Backpropagation and Predictive Coding]
Clearly, backpropagation is not local in all but the most simple examples. This is the reason why "`for most of the  past three decades since the invention of [backpropagation], it was generally believed that it could not be implemented by the brain"' \cite[][1837]{Liao2015}. In particular, this is the reason why a simple feed-forward network is most certainly not an implementation of backpropagation.\\
We will now be concerned with two things in this subsection: firstly, is the predictive coding model Hebbian? Secondly, is the predictive coding model an implementation of backpropagation. The first question will be answered in this paragraph and the second question will be answered by \thref{the:prediback}.\\
The learning rule on $\xi$, as given by \thref{learning-xi}, is clearly Hebbian. If we apply the learning rule to $\theta$ it is still local:
\[
\theta_{v\to \varepsilon^w}\mapsto \theta_{v\to\varepsilon^w}-\alpha N^{\varepsilon^w:i}(i)N^{v:o}(i),
\quad
\theta_{\varepsilon^w\to v}\mapsto\theta_{\varepsilon^w\to v}+\alpha N^{\varepsilon^w:i}(i)N^{v:o}(i).
\]
A challenge to its biological plausibility comes from the fact that the predictive coding model makes symmetric feedback and feed-forward weights necessary and "`there is no evidence for such a strong symmetry in cortex"' \cite[][1254]{Whittington2017}. However, Whittington and Bogacz remark that "`preliminary simulations suggest that symmetric weights are not necessary for good performance of predictive coding network"' \cite[][1254]{Whittington2017} and refer to a forthcoming paper. Interestingly, an investigation by Liao et al. \cite{Liao2015} provides similar results for a direct implementation of backpropagation.\\
On the other hand, the locality of the learning rule on $\lambda'$, as given by \thref{learning-lambda}, strongly depends on how $\Sigma_{\lambda'}$ is specified. In the case of $\Sigma_{\lambda'}=\lambda'$, though, it is clearly not local unless we only allow diagonal covariance matrices as $\Sigma$ needs to be inverted. Such a complicated computation poses a challenge for biological plausibility although Bogacz \cite{Bogacz2017} presents an alternative implementation with a local learning rule which we will investigate in chapter \ref{ch:further-work}.\\
In conclusion, while some problems still need to be reconciled, the predictive coding model is biologically more plausible than the backpropagation algorithm. This is not only supported by satisfying criteria like the Hebbian postulate but also by empirical evidence: the predictive coding model is based on a model presented by Rao and Ballard \cite{Rao1999} which they successfully employed to explain observations in visual processing neurons.
\end{Par}

\begin{Lem}\label{lem:predi-error-formula}
Consider some balanced state $s:=N(i)$ where $i\in S^{\mathcal{I}:i}[\mathfrak{P}],\mathcal{I}\subseteq\mathcal{G}$. For any $v\in\mathcal{G}$, if $i^v$ is not constant, 
\begin{align*}
s^{\varepsilon^v}=\left(f^v\right)'(s^{v:i})\xi_{v\to(v\to)}s^{\varepsilon^{v\to}}.
\end{align*}
\end{Lem}

\begin{Bew}
If $i^v$ is not constant, $s^v$ is not constant and therefore, the fixed point equation yields
\[
s^v=s^v-\alpha s^{\varepsilon^v}+\alpha\left(f^v\right)'(s^{v:i})\xi_{v\to}s^{\varepsilon}.
\]
For $w\notin v\to$, $\xi_{v\to w}=0$ which implies the lemma.
\end{Bew}

\begin{Par}
\thref{lem:predi-error-formula} is clearly equivalent to the error backpropagation in the feed-forward network $\mathfrak{N}$. However the states $s^{\mathcal{G}}$ may be different which is why the predictive coding network does not necessarily implement backpropagation. Note, however, that if $i\in S^{G^{(0)}:i}[\mathfrak{P}]$, $s^{\varepsilon^{\mathcal{G}\setminus G^{(0)}}}=\delta^{\mathcal{G}\setminus G^{(0)}}=0$ and the weight change is zero in both cases. On the other hand, in the standard case $i\in S^{G^{(0)}\cup G^{(L)}:i}[\mathfrak{P}]$, we do not necessarily obtain
\[
N[\mathfrak{P}](i)=N[\mathfrak{N}](i)
\]
and the weight change may be different. This is because the knowledge about $i^{G^{(L)}}$ provides further knowledge about the remaining PUs. It is thus intuitive that if the variance of $G^{(L)}$ is high, the remaining states are more similar to the states in the feed-forward network. This is supported by the following proposition:
\end{Par}

\begin{Pro}\label{the:prediback}
Consider a predictive coding model $\mathfrak{P}$ and a positive sequence $(a_n)_{n\in\mathbb{N}}$ such that $\lim_{n\to\infty}a_n=\infty$ and all $f_n$ are bounded. Consider a sequence of covariance matrices $(\Sigma^{(n)})_{n\in\mathbb{N}}$ where
\[
\Sigma^{(n)}=\begin{pmatrix}
\Sigma_{G^{(0)}\to G^{(0)}} & 0 & 0 \\
0 & \Sigma_{G^{(1,\dotsc, L-1)}\to G^{(1,\dotsc, L-1)}} & 0 \\
0 & 0 & a_n\cdot I_{G^{(L)}\to G^{(L)}}
\end{pmatrix}
\]
where $\Sigma_{G^{(0)}\to G^{(0)}}\in\mathbb{R}^{G^{(0)}\times G^{(0)}}$ and $\Sigma_{G^{(1,\dotsc, L-1)}\to G^{(1,\dotsc, L-1)}}\in\mathbb{R}^{G^{(L)}\times G^{(L)}}$ are some covariance matrices.\\
Let $i\in S^{G^{(0)}:i}[\mathfrak{P}],o\in S^{G^{(L)}:i}[\mathfrak{P}]$ be constant such that all $N_{\theta(\xi,\Sigma^{(n)}}[\mathfrak{P}](i,o)$ and $\lim_{n\to\infty}N_{\theta(\xi,\Sigma^{(n)}}[\mathfrak{P}](i,o)$ are well-defined. Then,
\begin{equation}\label{eq:convadd}
\lim_{n\to\infty} N_{\theta(\xi,\Sigma^{(n)})}^{\mathcal{G}}[\mathfrak{P}](i,o)=N_{\xi}[\mathfrak{N}](i)
\end{equation}
and, if 
\[
\tilde{\mathcal{L}}_{\alpha}:\mathbb{R}^{\mathcal{G}\times\mathcal{G}},
\quad
\tilde{\mathcal{L}}_{\alpha}^i(\xi,\Sigma):=\xi-\frac{\partial L_{\xi,\Sigma}}{\partial\xi},
\]
as given by \thref{learning-xi}, yields
\[
\mathcal{L}_{\alpha}:=\left(\tilde{\mathcal{L}}_{\alpha}\right)_{\const,\var:\mathbb{R}^{\mathcal{G}\times\mathcal{G}}},
\]
we obtain
\begin{equation}\label{eq:prediback}
\lim_{n\to\infty}\mathcal{L}_{\alpha a_n}^{(i,o)}(\xi,\Sigma^{(n)})=\mathcal{B}_{\alpha}^{i;o}(\xi).
\end{equation}
\end{Pro}

\begin{Bew}
Clearly,
\begin{align*}
\left(\Sigma^{(n)}\right)^{-1}=
\begin{pmatrix}
\Sigma_{G^{(0)}\to G^{(0)}}^{-1} & 0 & 0 \\
0 & \Sigma_{G^{(1,\dotsc, L-1)}\to G^{(1,\dotsc, L-1)}}^{-1} & 0 \\
0 & 0 & \frac{1}{a_n}\cdot I_{G^{(L)}\to G^{(L)}}
\end{pmatrix}.
\end{align*}
Defining
\[
s_n:=N_{\theta(\xi,\Sigma^{(n)})}[\mathfrak{P}](i,o),
\]
as $o\in S^{G^{(L):i}}[\mathfrak{P}]$ is constant, 
\begin{align*}
s_n^{\varepsilon^{G^{(L)}}}=\left(\Sigma_{G^{(L)}\to G^{(L)}}^{(n)}\right)^{-1}(s_n^{G^{(L)}:i}-\xi^T s_n^{\mathcal{G}:o})=\frac{1}{a_n}(o-\xi^Ts_n^{\mathcal{G}:o}).
\end{align*}
As $f_n$ are bounded, $(o-\xi^Ts_n^{\mathcal{G}:o})$ are bounded and, as $\lim_{n\to\infty}a_n=\infty$, 
\begin{align*}
\lim_{n\to\infty}s_n^{\varepsilon^{G^{(L)}}}=0.
\end{align*}
By induction, \thref{lem:predi-error-formula} implies
\[
\underset{v\in\mathcal{G}\setminus G^{(0)}}{\forall}s^{\varepsilon^v}=0.
\]
Analogously to \thref{lem:predicode-predistate} we can conclude (\ref{eq:convadd}). It is straight-forward to prove by induction that
\[
\lim_{n\to\infty}a_ns_n^{\varepsilon^{G^{(L)}}}=\lim_{n\to\infty}o-\xi^Ts_n^{\mathcal{G}:o}=o-\xi^TN_{\xi}[\mathfrak{N}](i)=o-N_{\xi}^{\mathcal{G}}[\mathfrak{N}](i)=\delta^{G^{(L)}}
\]
implies
\[
\lim_{n\to\infty}a_ns_n^{\varepsilon}=\delta
\]
as both formulas share a recursive, linear formula. Therefore,
\begin{align*}
\lim_{n\to\infty}\left(\mathcal{L}_{\alpha a_n}^{(i,o)}(\xi,\Sigma^{(n)})\right)_{v\to w}&=\xi_{v\to w}+\lim_{n\to\infty}\alpha a_ns_n^{\varepsilon^w}s_n^{v:o}=\\
\xi_{v\to w}+\lim_{n\to\infty}\delta^ws_n^{v:o}&=\xi_{v\to w}+\alpha \delta^w N^{v:o}_{\xi}[\mathfrak{N}](i)=\mathcal{B}_{\alpha}^{i;o}(\xi).
\end{align*}
\end{Bew}

\begin{Par}[Biological Plausibility III: Enforcing Backpropagation]
In this chapter, we have defined two kinds of networks: feed-forward networks $\mathfrak{N}$ with backpropagation learning and predictive coding networks $\mathfrak{P}$. Among the reasons for the definition of a predictive coding network was a biologically plausible implementation of backpropagation. \thref{the:prediback} yields such an implementation. The output variance simply has to be set high and be unlearnable and the resulting network satisfies locality of learning.\\
The relevance of such an implementation is, however, questionable. After all, predictive coding provides an entire model of reality, including hidden concepts whereas a feed-forward network only implements a blunt prediction of $y$ from $x$. The main reason why it is interesting to approximate backpropagation in a biologically plausible way is its proven effectiveness on the field of Machine Learning (see e. g. \cite{Whittington2017, Liao2015}). However, a similar argument may be made vice versa: predictive coding promises a flexible and meaningful statistical model. While its power remains to be proven, we may, once again be on the path of computational neuroscience inspiring statistics.\\
In summary, I would argue that biologically plausible implementations of backpropagation should not become an end in itself but rather launch a journey towards new models in neuroscience and statistics.\footnote{
I want to emphasize that I am also not suggesting that neuroscience faces such a risk of mathematization.} It would therefore be valuable to investigate the statistical properties of predictive coding as well as its (biologically plausible) capability to approximate other models. These questions motivate some initial investigations in the following chapter in which I hope to demonstrate that the mathematical potential of the framework I presented is far from exhausted.
\end{Par}

\chapter{Towards New Grounds: Open Problems}\label{ch:further-work}

\section{Normal Linear Model}

\section{General Linear Model}

\section{Network Pieces}

\chapter{Results and Outlook}

\printbibliography

\begin{appendices}
\chapter{Matrix Derivation}\label{ch:matrix-derivation}
\begin{Lem}\label{derivation-quadratic-form}
If $A\in\mathbb{R}^{n\times n},x\in\mathbb{R}^n$, $\frac{\partial x^TAx}{\partial x}=x^T(A+A^T)$.
\end{Lem}

\begin{Bew}
\begin{align*}
\frac{\partial}{\partial x_i}x^TAx=\frac{\partial}{\partial x_i}\sum_{j=1}^n\sum_{k=1}^nA_{jk}x_jx_k&=\sum_{j\ne i}A_{ji}x_i+\sum_{k\ne i}A_{ik}x_k+2A_{ii}x_i^2=\\
\sum_{j=1}^nA_{ji}x_i+\sum_{k=1}^nA_{ik}x_k&=x^T(A_{\cdot i}^T+A_{\cdot i}).
\end{align*}
\end{Bew}

\begin{Lem}\label{derivation-matrix-quadratic-form}
If $A\in\mathbb{R}^{n\times n},x\in\mathbb{R}^n$,
\begin{align*}
\frac{\partial x^TAx}{\partial A}=xx^T.
\end{align*}
\end{Lem}

\begin{Bew}
\begin{align*}
\frac{\partial}{\partial{A_{ij}}}x^TAx=\frac{\partial}{\partial{A_{i,j}}}\sum_{j=1}^n\sum_{k=1}^nA_jkx_jx_k=x_ix_j
\end{align*}
\end{Bew}

\begin{Lem}[Product Rule]\label{product-rule}
If
\[
M:\mathbb{R}\to\mathbb{R}^{l\times m},
\quad
N:\mathbb{R}\to\mathbb{R}^{m\times n},
\]
we obtain
\begin{align*}
\frac{\partial}{\partial x}(MN)(x)=M(x)\left(\frac{\partial}{\partial x}N(x)\right)+\left(\frac{\partial}{\partial x}M(x)\right)N(x).
\end{align*}
\end{Lem}

\begin{Bew}
\begin{align*}
\frac{\partial}{\partial x}(MN)_{ij}(x)=
\frac{\partial}{\partial x}M_i(x)N_j(x)=
\frac{\partial}{\partial x}\sum_{k=1}^mM_{ik}(x)N_{kj}(x)=\\
\sum_{k=1}^m
M_{ik}(x)\left(\frac{\partial}{\partial x}N_{kj}(x)\right)+
\left(\frac{\partial}{\partial x}M_{ik}(x)\right)N_{kj}(x)=
\left(M(x)\left(\frac{\partial}{\partial x}N(x)\right)\right)_{ij}+
\left(\left(\frac{\partial}{\partial x}M(x)\right)N(x)\right)_{ij}
\end{align*}
\end{Bew}

\begin{Lem}\label{inverse-chain}
\begin{align*}
\frac{\partial A(x)^{-1}}{\partial x}=-A(x)^{-1}\frac{\partial A(x)}{\partial x}A(x)^{-1}.
\end{align*}
\end{Lem}

\begin{Bew}
\begin{align*}
0=\frac{\partial}{\partial x}I=
\frac{\partial}{\partial x}A(x)A(x)^{-1}\overset{\thref{product-rule}}{=}
A(x)\left(\frac{\partial}{\partial x}A(x)^{-1}\right)+
\left(\frac{\partial}{\partial x}A(x)\right)A(x)^{-1},
\end{align*}
which implies the lemma.
\end{Bew}

\begin{Lem}[Chain Rule for matrix-valued functions]\label{chain-rule}
If
\[
M:\mathbb{R}\to\mathbb{R}^{m\times n}
\]
and
\[
f:\mathbb{R}^{m\times n}\to\mathbb{R}
\]
are partially differentiable, then
\[
f\circ M:\mathbb{R}\to\mathbb{R},
\quad
x\mapsto (f\circ M)(x)
\]
is differentiable and
\[
\frac{\partial f\left(M(x)\right)}{\partial x}=
\tr\left(\frac{\partial f\left(M(x)\right)}{\partial M(x)}\frac{\partial M(x)}{\partial x}\right).
\]
\end{Lem}

\begin{Bew}
We prove the lemma by transforming the matrix $M\in\mathbb{R}^{m\times n}$ into a vector $\tilde{M}\in\mathbb{R}^{mn}$ using the bijective mapping
\begin{align*}
\Psi:\mathbb{R}^{m\times n}\to\mathbb{R}^{mn},
\underset{\substack{i=1,\dotsc,m\\j=1,\dotsc,n}}{\forall}\Psi_{m(j-1)+i}(M):=M_{ij}.
\end{align*}
We define
\begin{align*}
\tilde{M}:=\Psi\circ M,
\quad
\tilde{f}:=f\circ\Psi^{-1},
\end{align*}
which implies
\begin{align*}
f\circ M=\tilde{f}\circ\tilde{M}.
\end{align*}
The matrix chain rule now follows directly from the chain rule for vector-valued functions although the notation is a bit cumbersome:
\begin{align*}
&\frac{\partial f\left(M(x)\right)}{\partial x}=
\frac{\partial\tilde{f}\left(\tilde{M}(x)\right)}{\partial x}=
\frac{\partial\tilde{f}\left(\tilde{M}(x)\right)}{\partial \tilde{M}(x)}\frac{\partial\tilde{M}(x)}{\partial x}=\\
&\sum_{i=1}^m\sum_{j=1}^n
\frac{\partial\tilde{f}\left(\tilde{M}(x)\right)}{\partial\tilde{M}_{m(j-1)+i}(x)}\cdot
\frac{\partial\tilde{M}_{m(j-1)+i}(x)}{\partial x}=
\sum_{i=1}^m\sum_{j=1}^n
\frac{\partial f\left(M(x)\right)}{\partial M_{ij}(x)}\cdot
\frac{\partial M_{ij}(x)}{\partial x}=\\
&\sum_{i=1}^m\sum_{j=1}^n
\left(\frac{\partial f\left(M(x)\right)}{\partial M(x)}\right)_{ji}\cdot
\left(\frac{\partial M_{ij}(x)}{\partial x}\right)_{ij}=
\sum_{j=1}^n\left(\frac{\partial f\left(M(x)\right)}{\partial M(x)}\frac{\partial M_{ij}(x)}{\partial x}\right)_{jj}=
\tr\left(\frac{\partial f\left(M(x)\right)}{\partial M(x)}\frac{\partial M(x)}{\partial x}\right).
\end{align*}
\end{Bew}

\begin{Lem}[Jacobi's Formula]\label{jacobis-formula}
If $\Lambda_i\subseteq\mathbb{R}$ is open and
\[
A:\Lambda_i\to\mathbb{R}^{n\times n}
\]
is a matrix-valued function, 
\[
\frac{\partial \det A(\lambda)}{\partial\lambda}=\tr\left((\text{Adj}A(\lambda)\frac{\partial}{\partial\lambda}A(\lambda)\right),
\]
where $\text{Adj}A(\lambda)$ is the classical adjoint matrix of $A(\lambda)$.\\
In particular, if $A$ is invertible,
\[
\frac{\partial \det A(\lambda)}{\partial\lambda}=\det A\tr\left(A(\lambda)^{-1}\frac{\partial}{\partial\lambda}A(\lambda)\right).
\]
\end{Lem}

\begin{Bew}
The proof can, for example, be found in \cite{Golberg1972}. The second equation is clear as the adjoint matrix suffices the formula
\begin{align*}
A\;\text{Adj(A)}=\det A\;I.
\end{align*}
\end{Bew}

\end{appendices}

\end{document}
